\documentclass[UTF8, 12pt, fontset=fandol, oneside]{ctexart}
\usepackage{researchnotes}

% 保留分册独立编译时的章号前缀。
\renewcommand{\thesection}{16.\arabic{section}}

\title{第十六章\quad 曲线积分与曲面积分}
\author{}
\date{}

\begin{document}

\maketitle

在 $\mathbb R^n(n\ge2)$ 中具有很多有趣的子集, 如 $\mathbb R^2$ 中的曲线, $\mathbb R^3$ 中的曲线与曲面等. 当 $n>3$ 时, 除了曲线与曲面外有趣的子集就更多了. 在科学技术的许多问题中, 人们会遇到这些子集上的积分问题. 在本章中, 我们将研究定义在 $\mathbb R^3$ (或 $\mathbb R^2$) 中一些特殊几何体上的积分问题, 主要考虑两种简单情况, 即曲线和曲面上的积分. 根据这些积分的不同背景, 人们将它们分成不同类型的积分. 我们主要介绍两类积分: 一类是多元函数的积分, 这类积分称为第一型积分; 另一类是向量函数的积分, 这类积分称为第二型积分. 在以下各节中, 我们将对它们作详细介绍。

\section{第一型曲线积分}

\subsection{第一型曲线积分的定义}

在定积分的应用中, 我们学习过曲线弧长的定义. 现在我们简要回顾一下. 设 $\Gamma$ 是平面内或空间 $\mathbb R^3$ 中的一条连续的简单曲线, 以 $A,B$ 为其端点 (当 $A=B$ 时, 则 $\Gamma$ 是一条约当曲线). 在 $\Gamma$ 上从 $A$ 到 $B$ 依次取点 $A=A_0,A_1,\cdots,A_K=B$, 它们构成了 $\Gamma$ 的一个分割 $T$. 记 $\overline{A_{k-1}A_k}$ 为连接 $A_{k-1},A_k$ 的线段, 其长为 $s_k(k=1,2,\cdots,K)$, 则当
\[
  \sup_T\sum_{k=1}^K s_k<+\infty
\]
时, 其中上确界是对 $\Gamma$ 所有的分割 $T$ 取的, 我们称该曲线是可求长的, 并称该上确界为 $\Gamma$ 的弧长。

现设光滑曲线 $\Gamma$ 由参数方程
\[
  \begin{cases}
    x=x(t),\\
    y=y(t),\\
    z=z(t)
  \end{cases}
  \quad(\alpha\le t\le \beta)
\]
给出, 则其弧长 $L$ 可由下述公式计算:
\[
  L=\int_\alpha^\beta \sqrt{(x'(t))^2+(y'(t))^2+(z'(t))^2}\mathrm dt.
\]

为了引入第一型曲线积分, 我们先来研究以下物质曲线的质量问题. 设 $\Gamma$ 是一条物质曲线, 它可被看成落在 $\mathbb R^3$ 中的一条曲线, 并在它上面定义了一个密度函数. 现在设 $\Gamma$ 以 $A$ 为起点, 以 $B$ 为终点, 在 $\Gamma$ 上任一点 $(x,y,z)$ 处该物质曲线的线密度为 $\rho(x,y,z)$, 我们来求 $\Gamma$ 的质量. 当线密度函数为常数函数, 即 $\rho(x,y,z)\equiv\rho_0$ 时, 我们可以求出该物质曲线的质量为 $\rho_0L$, 其中 $L$ 为该曲线的弧长. 当 $\rho(x,y,z)$ 在 $\Gamma$ 上不是常数函数时, 依照定积分的思想, 我们可以在 $\Gamma$ 上依次取点 $A=A_0,A_1,\cdots,A_K=B$, 并称其为 $\Gamma$ 的一个分割. 此分割将 $\Gamma$ 分成 $K$ 个小弧段. 在以 $A_{k-1},A_k$ 为端点的弧段 $\widehat{A_{k-1}A_k}$ 上任取一点 $(\xi_k,\eta_k,\zeta_k)(k=1,2,\cdots,K)$, 作和式 $\sum\limits_{k=1}^K \rho(\xi_k,\eta_k,\zeta_k)\Delta s_k$, 其中 $\Delta s_k(k=1,2,\cdots,K)$ 是 $\widehat{A_{k-1}A_k}$ 的弧长. 如同求曲边梯形的面积, 当分割越来越细时, 若该和式的极限趋向于常数 $I$, 则我们认为 $I$ 即为该物质曲线的质量. 由此我们引进下面的定义。

\noindent\textbf{定义 16.1.1}\quad 设 $\Gamma\subset\mathbb R^3$ 是可求长的简单曲线, 它以 $A$ 为起点, $B$ 为终点, $f(x,y,z)$ 是定义在 $\Gamma$ 上的函数. 对于 $\Gamma$ 的任一分割 $T$: $A=A_0,A_1,\cdots,A_K=B$, 记 $\Delta s_k(k=1,2,\cdots,K)$ 为弧段 $\widehat{A_{k-1}A_k}$ 的弧长 (即 $\Gamma$ 介于 $A_{k-1},A_k$ 之间部分的长度) 及 $\lambda(T)=\max\limits_{1\le k\le K}\Delta s_k$. 在 $\widehat{A_{k-1}A_k}$ 上任取一点 $(\xi_k,\eta_k,\zeta_k)(k=1,2,\cdots,K)$, 作和式
\[
  \sum_{k=1}^K f(\xi_k,\eta_k,\zeta_k)\Delta s_k.
\]
若存在 $I\in\mathbb R$, 使得对于 $\forall\varepsilon>0$, $\exists\delta>0$, 对 $\Gamma$ 的任意分割 $T$ 及每个小弧段上任意选取的 $(\xi_k,\eta_k,\zeta_k)$, 当 $\lambda(T)<\delta$ 时, 有
\[
  \left|\sum_{k=1}^K f(\xi_k,\eta_k,\zeta_k)\Delta s_k-I\right|<\varepsilon,
\]
即
\[
  \lim_{\lambda(T)\to0}\sum_{k=1}^K f(\xi_k,\eta_k,\zeta_k)\Delta s_k=I,
\]
则称 $f(x,y,z)$ 沿 $\Gamma$ 的第一型曲线积分存在, 并称 $I$ 为 $f(x,y,z)$ 在 $\Gamma$ 上的第一型曲线积分, 记为
\[
  I=\int_\Gamma f(x,y,z)\mathrm ds
  \quad\text{或}\quad
  I=\int_{\widehat{AB}} f(x,y,z)\mathrm ds,
\]
其中 $f(x,y,z)$ 称为被积函数, $\Gamma$ 称为积分曲线。

从定义 16.1.1 可以看出, 第一型曲线积分只与积分曲线 $\Gamma$ 的弧长及被积函数 $f(x,y,z)$ 的取值有关, 而与 $\Gamma$ 的方向无关. 设 $\Gamma=\widehat{AB}$, 即 $\Gamma$ 是以 $A$ 为起点, $B$ 为终点的一条曲线. 记 $\Gamma^-=\widehat{BA}$ 是 $\Gamma$ 的反向曲线, 即将 $\Gamma$ 的起点与终点对调后所得到的曲线, 则当 $f(x,y,z)$ 在 $\Gamma$ 上的第一型曲线积分存在时, 有
\[
  \int_\Gamma f(x,y,z)\mathrm ds=\int_{\Gamma^-} f(x,y,z)\mathrm ds.
\]

从定义 16.1.1 还可以推知第一型曲线积分具有下述性质:

(1) 若对于 $\forall(x,y,z)\in\Gamma$, 有 $f(x,y,z)\equiv1$, 则 $\int_\Gamma f(x,y,z)\mathrm ds=\int_\Gamma \mathrm ds$ 即为 $\Gamma$ 的弧长。

(2) 若 $\Gamma$ 落在 $Oxy$ 平面上, 且对于 $\forall(x,y)\in\Gamma$ 有 $f(x,y)>0$, 则 $\int_\Gamma f(x,y)\mathrm ds$ 表示空间柱面 $S=\{(x,y,z):(x,y)\in\Gamma,0\le z\le f(x,y)\}$ 的面积 (见图 16.1.1, 关于曲面面积的定义将在 §16.3 中讨论)。

(3) 第一型曲线积分具有关于被积函数的线性性, 即设 $\Gamma$ 是可求长的简单曲线, 函数 $f(x,y,z)$ 与 $g(x,y,z)$ 在 $\Gamma$ 上的第一型曲线积分存在, $\alpha,\beta\in\mathbb R$, 则 $\alpha f(x,y,z)+\beta g(x,y,z)$ 在 $\Gamma$ 上的第一型曲线积分存在, 并且
\[
  \int_\Gamma [\alpha f(x,y,z)+\beta g(x,y,z)]\mathrm ds
  =\alpha\int_\Gamma f(x,y,z)\mathrm ds
  +\beta\int_\Gamma g(x,y,z)\mathrm ds.
\]

(4) 第一型曲线积分具有关于积分曲线的可加性, 即设 $\Gamma_1=\widehat{AB}$, $\Gamma_2=\widehat{BC}$ 均为可求长的简单曲线, 且 $\Gamma=\Gamma_1+\Gamma_2=\widehat{AC}$ 也为可求长的简单曲线, 则函数 $f(x,y,z)$ 在 $\Gamma$ 上的第一型曲线积分存在的充分必要条件是 $f(x,y,z)$ 在 $\Gamma_1$ 和 $\Gamma_2$ 上的第一型曲线积分均存在, 且当第一型曲线积分存在时, 有
\[
  \int_{\Gamma_1} f(x,y,z)\mathrm ds+\int_{\Gamma_2} f(x,y,z)\mathrm ds
  =\int_\Gamma f(x,y,z)\mathrm ds.
\]
以上性质的证明请读者自己给出。

\subsection{第一型曲线积分的存在性与计算公式}

对于第一型曲线积分的存在性, 我们有下面的结论。

\noindent\textbf{定理 16.1.1}\quad 设 $\Gamma=\widehat{AB}\subset\mathbb R^3$ 是可求长的简单曲线, 函数 $f(x,y,z)$ 在 $\Gamma$ 上连续, 则第一型曲线积分 $\int_\Gamma f(x,y,z)\mathrm ds$ 存在。

\noindent\textbf{证明}\quad 设 $\Gamma$ 以弧长 $s$ 为参数的参数方程为
\[
  \Gamma:
  \begin{cases}
    x=x(s),\\
    y=y(s),\\
    z=z(s)
  \end{cases}
  \quad(0\le s\le L),
\]
其中 $L$ 为 $\Gamma$ 的弧长. 对于 $\Gamma=\widehat{AB}$ 的一个分割 $T$: $A=A_0,A_1,\cdots,A_K=B$, 则对应于闭区间 $[0,L]$ 的一个分割 $\Delta$: $0=s_0<s_1<\cdots<s_K=L$. 在 $\widehat{A_{k-1}A_k}$ 上任取一点 $(\xi_k,\eta_k,\zeta_k)$, 则必存在 $s'_k\in[s_{k-1},s_k]$, 使得 $(\xi_k,\eta_k,\zeta_k)=(x(s'_k),y(s'_k),z(s'_k))(k=1,2,\cdots,K)$. 因此
\[
  \sum_{k=1}^K f(\xi_k,\eta_k,\zeta_k)\Delta s_k
  =\sum_{k=1}^K f(x(s'_k),y(s'_k),z(s'_k))\Delta s_k.
\]
上述等式右端是连续函数 $f(x(s),y(s),z(s))$ 在区间 $[0,L]$ 上的黎曼和. 注意到 $\lambda(T)\to0$ 的充分必要条件是 $\lambda(\Delta)\to0$, 因此 $\int_\Gamma f(x,y,z)\mathrm ds$ 存在, 并且有
\[
  \int_\Gamma f(x,y,z)\mathrm ds=\int_0^L f(x(s),y(s),z(s))\mathrm ds.
\]
证毕。

在假定 $\Gamma\subset\mathbb R^3$ 是光滑曲线时, 我们有以下关于第一型曲线积分的计算公式。

\noindent\textbf{定理 16.1.2}\quad 设 $\Gamma$ 是光滑曲线, 其参数方程为
\[
  \begin{cases}
    x=x(t),\\
    y=y(t),\\
    z=z(t)
  \end{cases}
  \quad(\alpha\le t\le \beta),
\]
再设函数 $f(x,y,z)$ 在 $\Gamma$ 上连续, 则有
\[
  \int_\Gamma f(x,y,z)\mathrm ds
  =\int_\alpha^\beta f(x(t),y(t),z(t))
  \sqrt{[x'(t)]^2+[y'(t)]^2+[z'(t)]^2}\mathrm dt.
\]

\noindent\textbf{证明}\quad 由于 $\Gamma$ 是光滑曲线, 设其弧长为 $L$. 作 $[\alpha,\beta]\to[0,L]$ 的变换如下: 对于 $\forall t\in[\alpha,\beta]$, 记 $s(t)$ 为 $\Gamma$ 介于 $(x(\alpha),y(\alpha),z(\alpha))$ 与 $(x(t),y(t),z(t))$ 之间部分的弧长, 即
\[
  s=s(t)=\int_\alpha^t \sqrt{[x'(u)]^2+[y'(u)]^2+[z'(u)]^2}\mathrm du.
\]
注意到
\[
  \frac{\mathrm ds}{\mathrm dt}
  =\sqrt{[x'(t)]^2+[y'(t)]^2+[z'(t)]^2}>0,
\]
记 $t=t(s)$ 为 $s(t)$ 的反函数, 因此由定积分换元公式得
\[
\begin{aligned}
  \int_\Gamma f(x,y,z)\mathrm ds
  &=\int_0^L f(x(t(s)),y(t(s)),z(t(s)))\mathrm ds\\
  &=\int_\alpha^\beta f(x(t),y(t),z(t))
    \sqrt{[x'(t)]^2+[y'(t)]^2+[z'(t)]^2}\mathrm dt.
\end{aligned}
\]
证毕。

\noindent\textbf{例 16.1.1}\quad 计算第一型曲线积分 $\int_\Gamma \sqrt{1-x^2-y^2}\mathrm ds$, 其中 $\Gamma$ 是曲线 $x^2+y^2=x$。

\noindent\textbf{解法 1}\quad 取曲线 $\Gamma$ 的参数方程
\[
  \begin{cases}
    x=\dfrac12+\dfrac12\cos t,\\
    y=\dfrac12\sin t
  \end{cases}
  \quad(0\le t\le 2\pi),
\]
则
\[
  \mathrm ds=\sqrt{\frac14\sin^2t+\frac14\cos^2t}\mathrm dt
  =\frac12\mathrm dt.
\]
于是
\[
\begin{aligned}
  \int_\Gamma \sqrt{1-x^2-y^2}\mathrm ds
  &=\int_\Gamma \sqrt{1-x}\mathrm ds
    =\frac12\int_0^{2\pi}\sqrt{\frac{1-\cos t}{2}}\mathrm dt\\
  &=\frac12\int_0^{2\pi}\left|\sin\frac t2\right|\mathrm dt
    =\int_0^\pi\sin\frac t2\mathrm dt=2.
\end{aligned}
\]

\noindent\textbf{解法 2}\quad 由于被积函数及积分曲线的对称性, 我们只要考虑积分曲线落在第一象限的部分即可. 此时曲线 $\Gamma$ 可写成如下的参数方程
\[
  \begin{cases}
    x=x,\\
    y=\sqrt{x-x^2}
  \end{cases}
  \quad(0\le x\le1),
\]
从而 $\mathrm ds=\dfrac{\mathrm dx}{2\sqrt{x-x^2}}$, 因此
\[
  \int_\Gamma \sqrt{1-x^2-y^2}\mathrm ds
  =2\int_0^1\sqrt{1-x}\frac{\mathrm dx}{2\sqrt{x-x^2}}
  =\int_0^1\frac{\mathrm dx}{\sqrt x}=2.
\]

\noindent\textbf{注}\quad 在上例中, 解法 2 中的计算更简单些, 但最后的积分是一个瑕积分. 另外, 由第一型曲线积分的性质, 上例中的积分值是柱面 $x^2+y^2=x$ 被单位球面所截得的在 $Oxy$ 平面上方那一块的面积。

\noindent\textbf{例 16.1.2}\quad 计算第一型曲线积分 $\int_\Gamma (x+y)^2\mathrm ds$, 其中 $\Gamma$ 是单位圆盘上半部分的边界, 即 $\Gamma=\Gamma_1\cup\Gamma_2$, 其中
\[
  \Gamma_1=\{(x,y):x^2+y^2=1,y\ge0\},\quad
  \Gamma_2=\{(x,y):-1\le x\le1,y=0\}.
\]

\noindent\textbf{解}\quad 由第一型曲线积分的性质, 我们有
\[
  \int_\Gamma (x+y)^2\mathrm ds
  =\int_{\Gamma_1}(x+y)^2\mathrm ds+\int_{\Gamma_2}(x+y)^2\mathrm ds.
\]
由于
\[
  \int_{\Gamma_1}(x+y)^2\mathrm ds
  =\int_{\Gamma_1}(x^2+y^2)\mathrm ds+2\int_{\Gamma_1}xy\mathrm ds
  =\int_{\Gamma_1}\mathrm ds+0=\pi,
\]
其中 $\int_{\Gamma_1}xy\mathrm ds=0$ 由对称性得到, 另外我们有
\[
  \int_{\Gamma_2}(x+y)^2\mathrm ds
  =\int_{-1}^1 x^2\mathrm dx=\frac23,
\]
所以
\[
  \int_\Gamma (x+y)^2\mathrm ds=\frac23+\pi.
\]

\section{第二型曲线积分}

\subsection{第二型曲线积分的定义}

第二型曲线积分的物理模型是变力做功. 设空间有一个力场, 如地球对空间中任一点都存在引力, 该引力在空间的分布就可称为地球的引力场. 从数学的观点来看一个力场即是在 $\mathbb R^3$ 中定义了一个向量函数
\[
  \boldsymbol F(x,y,z)=(P(x,y,z),Q(x,y,z),R(x,y,z)).
\]
假定在它的作用下, 一个质点沿空间曲线 $\Gamma$ 从点 $A$ 移到了点 $B$. 现在我们来求该力场对质点所做的功。

首先注意到, 若 $\boldsymbol F$ 是常力, 即它的大小与方向均是常数, 并且在它的作用下, 质点沿连接点 $A$ 和 $B$ 的直线段 $\overline{AB}$ 从 $A$ 移到 $B$, 则它所做的功为
\[
  W=\boldsymbol F\cdot\overrightarrow{AB}.
\]
为了求解上述的一般问题, 我们用熟知的分割、求和、取极限的积分技巧. 记 $\Gamma=\widehat{AB}$, 沿 $\Gamma$ 依次取分点
\[
  A=A_0,A_1,\cdots,A_K=B.
\]
它们构成 $\Gamma$ 的一个分割 $T$, 记 $\lambda(T)=\max\limits_{1\le k\le K}\{\operatorname{diam}(\widehat{A_{k-1}A_k})\}$. 在每一段小弧 $\widehat{A_{k-1}A_k}(k=1,2,\cdots,K)$ 上, 以线段 $\overline{A_{k-1}A_k}$ 代替 $\widehat{A_{k-1}A_k}$, 并且将 $\boldsymbol F(x,y,z)$ 在 $\widehat{A_{k-1}A_k}$ 上取常力 $\boldsymbol F(\xi_k,\eta_k,\zeta_k)((\xi_k,\eta_k,\zeta_k)\in\widehat{A_{k-1}A_k})$, 则在变力 $\boldsymbol F$ 作用下质点沿 $\Gamma$ 从 $A_{k-1}$ 移动到 $A_k$ 所做的功
\[
  \Delta W_k\approx \boldsymbol F(\xi_k,\eta_k,\zeta_k)\cdot\overrightarrow{A_{k-1}A_k}.
\]
因此
\[
  W=\sum_{k=1}^K\Delta W_k
  \approx \sum_{k=1}^K \boldsymbol F(\xi_k,\eta_k,\zeta_k)\cdot\overrightarrow{A_{k-1}A_k}.
\]
当 $\lambda(T)\to0$ 时, 若上式右端和式的极限存在, 则
\[
  W=\lim_{\lambda(T)\to0}\sum_{k=1}^K
  \boldsymbol F(\xi_k,\eta_k,\zeta_k)\cdot\overrightarrow{A_{k-1}A_k}.
\]
记 $A_k=(x_k,y_k,z_k)(k=1,2,\cdots,K)$, 则
\[
  \overrightarrow{A_{k-1}A_k}
  =(x_k-x_{k-1},y_k-y_{k-1},z_k-z_{k-1})
  \triangleq(\Delta x_k,\Delta y_k,\Delta z_k).
\]
于是我们有
\[
\begin{aligned}
  &\sum_{k=1}^K \boldsymbol F(\xi_k,\eta_k,\zeta_k)\cdot\overrightarrow{A_{k-1}A_k}\\
  &\quad=\sum_{k=1}^K\bigl(P(\xi_k,\eta_k,\zeta_k)\Delta x_k
    +Q(\xi_k,\eta_k,\zeta_k)\Delta y_k
    +R(\xi_k,\eta_k,\zeta_k)\Delta z_k\bigr).
\end{aligned}
\]
从上述变力做功的问题中, 我们得到一个新的和式极限. 由此和式极限我们可以引进一类新的积分. 另外要注意, 我们记空间一条曲线 $\Gamma=\widehat{AB}$ 是为了强调 $\Gamma$ 是以 $A$ 为起点, $B$ 为终点的. 当 $A=B$ 时, 我们则必须在 $\Gamma$ 上取定一个方向。

\noindent\textbf{定义 16.2.1}\quad 设 $\Gamma=\widehat{AB}$ 为空间连续曲线, 向量函数
\[
  \boldsymbol F(x,y,z)=(P(x,y,z),Q(x,y,z),R(x,y,z))
\]
在 $\widehat{AB}$ 上有定义, 记在 $\widehat{AB}$ 上依次取分点
\[
  A=A_0,A_1,\cdots,A_K=B
\]
组成的分割为 $T$,\allowbreak $\overrightarrow{A_{k-1}A_k}=(\Delta x_k,\Delta y_k,\Delta z_k)(k=1,2,\cdots,K)$, $\lambda(T)=\max\limits_{1\le k\le K}\{\operatorname{diam}(\widehat{A_{k-1}A_k})\}$. 在 $\widehat{A_{k-1}A_k}$ 上任取点 $(\xi_k,\eta_k,\zeta_k)(k=1,2,\cdots,K)$. 若极限
\[
  \lim_{\lambda(T)\to0}\sum_{k=1}^K
  [P(\xi_k,\eta_k,\zeta_k)\Delta x_k
  +Q(\xi_k,\eta_k,\zeta_k)\Delta y_k
  +R(\xi_k,\eta_k,\zeta_k)\Delta z_k]
\]
存在并且等于 $I$, 其中常数 $I$ 不依赖 $\widehat{AB}$ 的分割 $T$ 以及 $(\xi_k,\eta_k,\zeta_k)$ 的选取, 则称该极限值 $I$ 为 $\boldsymbol F(x,y,z)$ 在 $\Gamma$ 上的第二型曲线积分, 记为
\[
  I=\int_\Gamma P(x,y,z)\mathrm dx+Q(x,y,z)\mathrm dy+R(x,y,z)\mathrm dz,
\]
或者
\[
  I=\int_{\widehat{AB}}P\mathrm dx+Q\mathrm dy+R\mathrm dz,
\]
其中 $\boldsymbol F(x,y,z)=(P,Q,R)$ 称为被积函数, $\Gamma$ 称为积分曲线。

从上述定义可以看出, 第二型曲线积分是考虑一个向量函数在曲线上沿着一个方向的积分, 因此它与曲线的方向有关. 显然, 第一型曲线积分是考虑一个多元函数在曲线上的积分. 当然, 仅从纯数学的观点, 第二型曲线积分也可仅仅定义一个多元函数 $P(x,y,z)$ 的下述积分 $\int_{\widehat{AB}}P\mathrm dx$, $\int_{\widehat{AB}}P\mathrm dy$, $\int_{\widehat{AB}}P\mathrm dz$. 但是这些积分不能反映出第二型曲线积分的物理背景。

从定义我们可以看出第二型曲线积分具有下述性质:

(1) 若 $\int_{\widehat{AB}}P\mathrm dx+Q\mathrm dy+R\mathrm dz$ 存在, 则 $\int_{\widehat{BA}}P\mathrm dx+Q\mathrm dy+R\mathrm dz$ 存在, 并且
\[
  \int_{\widehat{BA}}P\mathrm dx+Q\mathrm dy+R\mathrm dz
  =-\int_{\widehat{AB}}P\mathrm dx+Q\mathrm dy+R\mathrm dz.
\]
特别地, 当 $\Gamma$ 是一条简单闭曲线时, 我们有
\[
  \int_{\Gamma^-}P\mathrm dx+Q\mathrm dy+R\mathrm dz
  =-\int_\Gamma P\mathrm dx+Q\mathrm dy+R\mathrm dz.
\]

(2) 设向量函数 $\boldsymbol f=(P_1,Q_1,R_1)$ 和 $\boldsymbol g=(P_2,Q_2,R_2)$ 在空间曲线 $\widehat{AB}$ 上的第二型曲线积分都存在, $\alpha,\beta\in\mathbb R$, 则 $\alpha\boldsymbol f+\beta\boldsymbol g$ 在 $\widehat{AB}$ 上的第二型曲线积分存在, 并且
\[
\begin{aligned}
  &\int_{\widehat{AB}}(\alpha P_1\mathrm dx+\alpha Q_1\mathrm dy+\alpha R_1\mathrm dz)
  +(\beta P_2\mathrm dx+\beta Q_2\mathrm dy+\beta R_2\mathrm dz)\\
  &\quad=\alpha\int_{\widehat{AB}}P_1\mathrm dx+Q_1\mathrm dy+R_1\mathrm dz
  +\beta\int_{\widehat{AB}}P_2\mathrm dx+Q_2\mathrm dy+R_2\mathrm dz.
\end{aligned}
\]

(3) 设 $\widehat{AB}=\widehat{AC}\cup\widehat{CB}$, 即 $\widehat{AB}$ 是由两条连续曲线 $\widehat{AC},\widehat{CB}$ 组成, 并且 $\widehat{AC}$ 与 $\widehat{CB}$ 仅在端点 $C$ 处相交, 则向量函数 $\boldsymbol F=(P,Q,R)$ 在 $\widehat{AB}$ 上的第二型曲线积分存在的充分必要条件是它在 $\widehat{AC}$ 及 $\widehat{CB}$ 上的第二型曲线积分均存在, 并且第二型曲线积分存在时, 成立等式
\[
  \int_{\widehat{AB}}P\mathrm dx+Q\mathrm dy+R\mathrm dz
  =\int_{\widehat{AC}}P\mathrm dx+Q\mathrm dy+R\mathrm dz
  +\int_{\widehat{CB}}P\mathrm dx+Q\mathrm dy+R\mathrm dz.
\]

\subsection{第二型曲线积分的存在性与计算公式}

在这里, 我们不准备讨论关于非常广泛的积分曲线和被积函数的第二型曲线积分存在性问题, 只讨论积分曲线是光滑曲线, 而被积函数是连续的情形。

\noindent\textbf{定理 16.2.1}\quad 设 $\Gamma=\widehat{AB}$ 是以 $A$ 为起点, $B$ 为终点的光滑曲线, 其参数方程为
\[
  \begin{cases}
    x=x(t),\\
    y=y(t),\\
    z=z(t)
  \end{cases}
  \quad(\alpha\le t\le \beta)
\]
且当 $t$ 从 $\alpha$ 连续变为 $\beta$ 时, 对应曲线上的点从 $A$ 连续变为 $B$, 再假定函数\allowbreak $P(x,y,z),Q(x,y,z),R(x,y,z)$ 在 $\Gamma$ 上连续, 则向量函数 $\boldsymbol F=(P,Q,R)$ 在 $\widehat{AB}$ 上的第二型曲线积分存在, 并且
\[
\begin{aligned}
  \int_{\widehat{AB}}P\mathrm dx+Q\mathrm dy+R\mathrm dz
  &=\int_\alpha^\beta P(x(t),y(t),z(t))x'(t)\mathrm dt
    +\int_\alpha^\beta Q(x(t),y(t),z(t))y'(t)\mathrm dt\\
  &\quad+\int_\alpha^\beta R(x(t),y(t),z(t))z'(t)\mathrm dt.
\end{aligned}
\]

\noindent\textbf{证明}\quad 我们用微元法证之. 由光滑曲线的定义知 $x'(t),y'(t),z'(t)$ 均是 $t$ 的连续函数, 且对于 $\forall t\in(\alpha,\beta)$, 有
\[
  x'^2(t)+y'^2(t)+z'^2(t)\ne0.
\]
对于 $\forall[t,t+\mathrm dt]\subset[\alpha,\beta]$, 记 $A_t=(x(t),y(t),z(t))\in\Gamma$, $A_{t+\mathrm dt}=(x(t+\mathrm dt),y(t+\mathrm dt),z(t+\mathrm dt))\in\Gamma$, 由于
\[
\begin{aligned}
  \overrightarrow{A_tA_{t+\mathrm dt}}
  &\approx (x'(t),y'(t),z'(t))\mathrm dt\\
  &=\frac{(x'(t),y'(t),z'(t))}
    {\sqrt{x'^2(t)+y'^2(t)+z'^2(t)}}
    \sqrt{x'^2(t)+y'^2(t)+z'^2(t)}\mathrm dt\\
  &=\frac{(x'(t),y'(t),z'(t))}
    {\sqrt{x'^2(t)+y'^2(t)+z'^2(t)}}\mathrm ds\\
  &=(\cos\alpha,\cos\beta,\cos\gamma)\mathrm ds.
\end{aligned}
\]
其中 $(\cos\alpha,\cos\beta,\cos\gamma)$ 是曲线在点 $(x(t),y(t),z(t))$ 处的单位切向量. 因此有
\[
\begin{aligned}
  P\mathrm dx+Q\mathrm dy+R\mathrm dz
  &=(P,Q,R)(\mathrm dx,\mathrm dy,\mathrm dz)\\
  &=(P,Q,R)(\cos\alpha,\cos\beta,\cos\gamma)\mathrm ds\\
  &=(P\cos\alpha+Q\cos\beta+R\cos\gamma)\mathrm ds.
\end{aligned}
\]
由假设, $\Gamma$ 是光滑曲线, 从而 $\cos\alpha,\cos\beta,\cos\gamma$ 是 $t$ 的连续函数. 由第一型曲线积分的存在性定理知 $\int_{\widehat{AB}}(P\cos\alpha+Q\cos\beta+R\cos\gamma)\mathrm ds$ 存在, 从而 $\int_{\widehat{AB}}P\mathrm dx+Q\mathrm dy+R\mathrm dz$ 存在, 并且有公式
\[
\begin{aligned}
  \int_{\widehat{AB}}P\mathrm dx+Q\mathrm dy+R\mathrm dz
  &=\int_\alpha^\beta P(x(t),y(t),z(t))x'(t)\mathrm dt
    +\int_\alpha^\beta Q(x(t),y(t),z(t))y'(t)\mathrm dt\\
  &\quad+\int_\alpha^\beta R(x(t),y(t),z(t))z'(t)\mathrm dt.
\end{aligned}
\]
证毕。

\noindent\textbf{注}\quad 在上面定理的证明过程中, 我们清楚地知道, 第二型曲线积分与第一型曲线积分有如下的关系:
\[
  \int_{\widehat{AB}}P\mathrm dx+Q\mathrm dy+R\mathrm dz
  =\int_{\widehat{AB}}(P,Q,R)\cdot(\cos\alpha,\cos\beta,\cos\gamma)\mathrm ds,
\]
其中 $(\cos\alpha,\cos\beta,\cos\gamma)$ 为 $\widehat{AB}$ 在 $(x,y,z)$ 处的单位切向量。

到现在为止, 我们讨论了平面及空间的曲线积分问题. 值得指出的是, 在 $\mathbb R^n(n>3)$ 中也可以讨论相应的曲线积分问题. 事实上, 我们已经有了 $\mathbb R^n$ 中曲线的定义, 因此可以用完全平行于 $\mathbb R^3$ 中的理论来讨论可求长曲线、光滑曲线等问题, 从而可以定义两类不同的积分. 由于这些理论与我们已经研究的 $\mathbb R^3$ 中的相应理论有很大的相似性, 我们在这里就不深入讨论了。

\noindent\textbf{例 16.2.1}\quad 计算第二型曲线积分 $\int_{\widehat{AB}}(-y)\mathrm dx+x\mathrm dy$, 其中 $\widehat{AB}$ 为单位圆周 $x^2+y^2=1$ 的上半部分, 方向为从点 $A(1,0)$ 到点 $B(-1,0)$。

\noindent\textbf{解法 1}\quad 令
\[
  \begin{cases}
    x=\cos t,\\
    y=\sin t
  \end{cases}
  \quad(0\le t\le\pi),
\]
则 $t$ 从 $0$ 连续变化到 $\pi$ 时, 曲线上的点从 $A$ 连续变化到 $B$. 因此
\[
\begin{aligned}
  \int_{\widehat{AB}}(-y)\mathrm dx+x\mathrm dy
  &=\int_0^\pi[-\sin t\cdot(-\sin t)+\cos t\cdot\cos t]\mathrm dt\\
  &=\int_0^\pi\mathrm dt=\pi.
\end{aligned}
\]

\noindent\textbf{解法 2}\quad 由 $x^2+y^2=1$ 知 $x\mathrm dx+y\mathrm dy=0$, 得 $\mathrm dy=-\dfrac{x}{y}\mathrm dx$, 从而有
\[
\begin{aligned}
  \int_{\widehat{AB}}(-y)\mathrm dx+x\mathrm dy
  &=\int_1^{-1}(-y)\mathrm dx+\int_1^{-1}x\left(-\frac{x}{y}\right)\mathrm dx\\
  &=\int_{-1}^1\left(\frac{x^2+y^2}{y}\right)\mathrm dx
    =\int_{-1}^1\frac{\mathrm dx}{\sqrt{1-x^2}}
    =\arcsin x\bigg|_{-1}^1=\pi.
\end{aligned}
\]

\noindent\textbf{例 16.2.2}\quad 计算第二型曲线积分 $\int_\Gamma (x^2-y^2)\mathrm dx-2xy\mathrm dy$, 其中 $\Gamma$ 是从点 $(0,0)$ 沿曲线 $y=x^\alpha(\alpha>0)$ 到点 $(1,1)$ 的部分。

\noindent\textbf{解}\quad 将曲线
\[
  \begin{cases}
    x=x,\\
    y=x^\alpha
  \end{cases}
  \quad(x\in[0,1])
\]
代入得
\[
\begin{aligned}
  \int_\Gamma (x^2-y^2)\mathrm dx-2xy\mathrm dy
  &=\int_0^1(x^2-x^{2\alpha})\mathrm dx
    -\int_0^1 2xx^\alpha(\alpha x^{\alpha-1})\mathrm dx\\
  &=\int_0^1[x^2-(2\alpha+1)x^{2\alpha}]\mathrm dx
    =-\frac23.
\end{aligned}
\]
此题说明, 该第二型曲线积分与 $\alpha$ 的选取无关, 而只与这些曲线的端点有关。

\noindent\textbf{例 16.2.3}\quad 计算第二型曲线积分 $\int_\Gamma x\mathrm dx+y\mathrm dy+z\mathrm dz$, 其中 $\Gamma$ 为球面 $x^2+y^2+z^2=1$ 与平面 $x+y+z=0$ 的交线, 从 $z$ 轴看去取逆时针方向。

\noindent\textbf{解法 1}\quad 作正交变换将 $\Gamma$ 变到 $O\xi\eta$ 平面:
\[
  \begin{cases}
    \xi=\dfrac1{\sqrt2}(x-y),\\
    \eta=\dfrac1{\sqrt6}(x+y-2z),\\
    \zeta=\dfrac1{\sqrt3}(x+y+z).
  \end{cases}
\]
然后利用 $\xi=\cos t,\eta=\sin t$ 代入上述 $\Gamma$ 的方程, 即得 $\Gamma$ 的参数方程
\[
  \begin{cases}
    x=\dfrac1{\sqrt2}\cos t+\dfrac1{\sqrt6}\sin t,\\
    y=-\dfrac1{\sqrt2}\cos t+\dfrac1{\sqrt6}\sin t,\\
    z=-\dfrac2{\sqrt6}\sin t,
  \end{cases}
\]
从而
\[
\begin{aligned}
  \int_\Gamma x\mathrm dx+y\mathrm dy+z\mathrm dz
  &=\int_0^{2\pi}\bigg[
    \left(\frac1{\sqrt2}\cos t+\frac1{\sqrt6}\sin t\right)
    \left(-\frac1{\sqrt2}\sin t+\frac1{\sqrt6}\cos t\right)\\
  &\quad+
    \left(-\frac1{\sqrt2}\cos t+\frac1{\sqrt6}\sin t\right)
    \left(\frac1{\sqrt2}\sin t+\frac1{\sqrt6}\cos t\right)\\
  &\quad+
    \left(-\frac2{\sqrt6}\sin t\right)
    \left(-\frac2{\sqrt6}\cos t\right)
  \bigg]\mathrm dt\\
  &=\int_0^{2\pi}0\mathrm dt=0.
\end{aligned}
\]

\noindent\textbf{解法 2}\quad 由于球面 $x^2+y^2+z^2=1$ 上每个点 $(x,y,z)$ 在 $\mathbb R^3$ 中表示的向量即是球面在该点的单位法向量 $\boldsymbol n$, 而曲线 $\Gamma$ 每个点处的单位切向量 $\boldsymbol v$ 与球面在该点的法向量正交, 因此有
\[
  \int_\Gamma x\mathrm dx+y\mathrm dy+z\mathrm dz
  =\int_\Gamma \boldsymbol n\cdot\boldsymbol v\mathrm ds=0.
\]

\section{第一型曲面积分}

\subsection{曲面的面积}

在介绍第一型曲面积分之前, 我们先来讨论一下曲面的面积问题. 如何利用积分的思想来求曲面的面积呢? 受求曲线弧长的启发, 我们能否用曲面内接多边形的面积来逼近曲面的面积? 对此施瓦茨曾经有一个著名的例子, 他考虑了由圆柱面的内接全等等腰三角形构成的多面体面积的逼近: 设圆柱面 $S$ 为 $\{(x,y,z):x^2+y^2=1,0\le z\le1\}$. 将 $S$ 的高进行了 $K$ 等分, 从而在 $S$ 上得到 $K+1$ 个平行于 $Oxy$ 平面的圆周 $\Gamma_k(k=1,2,\cdots,K+1)$. 对每个圆周 $\Gamma_k(k=1,2,\cdots,K+1)$ 再 $L$ 等分, 使得 $\Gamma_{k+1}$ 上的分点在 $Oxy$ 平面的投影落在 $\Gamma_k$ 上的分点在 $Oxy$ 平面的投影的中间, 然后将每个 $\Gamma_k$ 上相邻两个分点连接, 并且它们分别与 $\Gamma_{k+1}$ 上的位于它上方中间的点连接, 这样我们就得到了一个等腰三角形. 所有的这些三角形组成了 $S$ 的一个内接多边形. 这个多边形的面积用初等数学的知识即可求出. 容易证明, 适当选取 $K,L$ 的比例, 如 $K=L^3$, 则可以使得多边形的面积趋于 $\infty$. 在上述逼近中, 若取 $K=L$, 则容易推出多边形的面积趋于 $2\pi$. 因此用内接多边形的面积来逼近曲面面积的方法是行不通的。

仔细分析一下施瓦茨的例子, 当 $K$ 与 $L$ 的比很大时, 即三角形的高与底边长的比趋于零时, 每个小三角形所在平面几乎与柱面的切平面垂直, 因此产生很多由于弯曲而增加的面积. 如果该三角形所在平面越来越“平行”于切平面, 则多边形面积逼近将会成功. 注意到对于圆柱面, 我们若将其摊平到平面, 则可以计算出它的面积为
\[
  \text{底周长}\times\text{高}=2\pi.
\]

为了求曲面面积, 我们的想法是用每一点处切平面上的部分来构造曲面的外切多边形进行逼近. 由于要用到切平面, 我们只考虑光滑曲面。

设曲面 $S$ 由参数方程
\begin{equation}
  \begin{cases}
    x=x(u,v),\\
    y=y(u,v),\\
    z=z(u,v),
  \end{cases}
  \quad (u,v)\in D
  \tag{16.3.1}
\end{equation}
定义, 其中 $D$ 是可求面积的有界闭区域, 函数 $x(u,v),y(u,v),z(u,v)$ 均在 $D$ 上具有连续偏导数. 作为 $D$ 到 $\mathbb R^3$ 的映射, (16.3.1) 是单的且其雅可比矩阵
\[
  \begin{pmatrix}
    x'_u & y'_u & z'_u\\
    x'_v & y'_v & z'_v
  \end{pmatrix}
\]
在每一点 $(u,v)$ 处的秩均为 $2$. 在本节中, 我们总假定曲面 $S$ 为满足上述条件的光滑曲面。

现记 $\boldsymbol r(u,v)=(x(u,v),y(u,v),z(u,v))$ 及
\begin{equation}
  A=\frac{\partial(y,z)}{\partial(u,v)},\quad
  B=\frac{\partial(z,x)}{\partial(u,v)},\quad
  C=\frac{\partial(x,y)}{\partial(u,v)}.
  \tag{16.3.2}
\end{equation}
则在任一点 $(u,v)$ 处行列式 $A,B,C$ 中至少有一个不为零, 并且它们均是 $(u,v)$ 的连续函数, 因此法向量 $\boldsymbol n=(A,B,C)$ 在曲面 $S$ 上连续变化。

下面我们利用微元法来求曲面的面积. 给定以 $(u_0,v_0),(u_0+\mathrm du,v_0),(u_0+\mathrm du,v_0+\mathrm dv),(u_0,v_0+\mathrm dv)$ 为顶点的一个小矩形 $D_0$. 当 $(u,v)$ 在 $D_0$ 上变化时, 在 $S$ 上可得到一小块曲面, 它可近似看成一个小平行四边形, 其两邻边分别为
\[
  \boldsymbol r(u_0+\mathrm du,v_0)-\boldsymbol r(u_0,v_0)
  \approx \boldsymbol r'_u(u_0,v_0)\mathrm du
\]
与
\[
  \boldsymbol r(u_0,v_0+\mathrm dv)-\boldsymbol r(u_0,v_0)
  \approx \boldsymbol r'_v(u_0,v_0)\mathrm dv.
\]
因此它们张成的平行四边形的面积, 即面积微元为
\[
  \mathrm dS(u_0,v_0)=|\boldsymbol r'_u(u_0,v_0)\times\boldsymbol r'_v(u_0,v_0)|\mathrm du\mathrm dv,
\]
从而曲面 $S$ 的面积 (仍用 $S$ 表示其面积) 为
\[
  S=\iint_D |\boldsymbol r'_u(u,v)\times\boldsymbol r'_v(u,v)|\mathrm du\mathrm dv.
\]
现在我们进一步计算 $|\boldsymbol r'_u(u,v)\times\boldsymbol r'_v(u,v)|$. 由
\[
  |\boldsymbol r'_u(u,v)\times\boldsymbol r'_v(u,v)|
  =|(x'_u,y'_u,z'_u)\times(x'_v,y'_v,z'_v)|
  =\sqrt{A^2+B^2+C^2},
\]
若记
\begin{equation}
  \begin{cases}
    E=\boldsymbol r'_u\boldsymbol r'_u=x_u'^2+y_u'^2+z_u'^2,\\
    F=\boldsymbol r'_u\boldsymbol r'_v=x'_ux'_v+y'_uy'_v+z'_uz'_v,\\
    G=\boldsymbol r'_v\boldsymbol r'_v=x_v'^2+y_v'^2+z_v'^2,
  \end{cases}
  \tag{16.3.3}
\end{equation}
则可知
\[
  EG-F^2=A^2+B^2+C^2.
\]
因此, 我们有下面的定理。

\noindent\textbf{定理 16.3.1}\quad 设光滑曲面 $S$ 由参数方程 (16.3.1) 给出, 则它的面积为
\[
  S=\iint_D \sqrt{A^2+B^2+C^2}\mathrm du\mathrm dv
  =\iint_D \sqrt{EG-F^2}\mathrm du\mathrm dv,
\]
其中 $A,B,C$ 和 $E,F,G$ 的定义分别见式 (16.3.2) 和 (16.3.3)。

特别地, 当光滑曲面 $S$ 是由 $z=f(x,y)((x,y)\in D)$ 给出时, 则有
\[
  S=\iint_D \sqrt{1+f_x'^2(x,y)+f_y'^2(x,y)}\mathrm dx\mathrm dy.
\]

\noindent\textbf{例 16.3.1}\quad 求单位圆柱面 $x^2+y^2=1$ 在平面 $z=0$ 与 $z=1$ 之间的面积。

\noindent\textbf{解}\quad 设所求面积为 $S$, 则 $S$ 为曲面 $y=f(z,x)=\sqrt{1-x^2}((z,x)\in D,\text{ 其中 }D=[0,1]\times[-1,1])$ 的面积的两倍. 因此有
\[
\begin{aligned}
  S&=2\iint_D\sqrt{1+f_x'^2+f_z'^2}\mathrm dz\mathrm dx\\
  &=2\int_0^1\mathrm dz\int_{-1}^1\sqrt{1+\left(\frac{-x}{\sqrt{1-x^2}}\right)^2}\mathrm dx\\
  &=2\int_{-1}^1\frac{\mathrm dx}{\sqrt{1-x^2}}=2\pi.
\end{aligned}
\]

\noindent\textbf{例 16.3.2}\quad 计算球面 $x^2+y^2+z^2=1$ 被柱面 $\left(x-\dfrac12\right)^2+y^2=\dfrac14$ 所截在柱面内的部分的面积。

\noindent\textbf{解}\quad 由对称性, 所求面积 $S$ 为所截得曲面在 $Oxy$ 平面上方部分面积的两倍. 设 $z=f(x,y)=\sqrt{1-x^2-y^2}((x,y)\in D)$, 其中
\[
  D=\left\{(x,y):\left(x-\frac12\right)^2+y^2\le\frac14\right\},
\]
则
\[
  S=2\iint_D\sqrt{1+f_x'^2+f_y'^2}\mathrm dx\mathrm dy
  =2\iint_D\frac{\mathrm dx\mathrm dy}{\sqrt{1-x^2-y^2}}.
\]
利用极坐标变换得
\[
\begin{aligned}
  S&=4\int_0^{\frac\pi2}\mathrm d\theta\int_0^{\cos\theta}
  \frac{r\mathrm dr}{\sqrt{1-r^2}}
  =4\int_0^{\frac\pi2}\left(-\sqrt{1-r^2}\right)\bigg|_0^{\cos\theta}\mathrm d\theta\\
  &=4\int_0^{\frac\pi2}(1-\sin\theta)\mathrm d\theta=2\pi-4.
\end{aligned}
\]

\subsection{第一型曲面积分的定义}

由于曲面的复杂性, 我们在这里只讨论光滑曲面上的第一型曲面积分. 具体来说, 在本小节中我们所指的曲面 $S$ 将由参数方程 (16.3.1) 定义, 并且它们满足所假定的光滑条件. 特别地, 我们总假定闭区域 $D\subset\mathbb R^2$ 的边界由有限条光滑曲线组成. 设
\[
  \Delta=\{\Delta D_1,\Delta D_2,\cdots,\Delta D_K\}
\]
为 $D$ 的一个分割, 且假定每个 $\Delta D_k(k=1,2,\cdots,K)$ 均为由有限条光滑曲线所围成的区域. 在此假定下, $\Delta S_k=\{\boldsymbol r(u,v):(u,v)\in\Delta D_k\}(k=1,2,\cdots,K)$ 是一块光滑小曲面, $T=\{\Delta S_1,\Delta S_2,\cdots,\Delta S_K\}$ 是 $S$ 的一个分割. 记 $\lambda(\Delta)=\max\limits_{1\le k\le K}\operatorname{diam}(\Delta D_k)$ 与 $\lambda(T)=\max\limits_{1\le k\le K}\operatorname{diam}(\Delta S_k)$. 由于 $\boldsymbol r(u,v)$ 在 $D$ 上一致连续且其逆映射在 $S$ 上一致连续, 因此 $\lambda(\Delta)\to0$ 的充分必要条件是 $\lambda(T)\to0$。

对于 $S$ 的分割 $T=\{\Delta S_1,\Delta S_2,\cdots,\Delta S_K\}$, 以下我们总是假定每个 $\Delta S_k$ 均是 $S$ 上可求面积的一小块曲面, 并且存在 $D$ 的可求面积的闭子区域 $\Delta D_k$, 使得 $\boldsymbol r(u,v)$ 是 $\Delta D_k$ 到 $\Delta S_k$ 的一一对应, 从而 $\Delta=\{\Delta D_1,\Delta D_2,\cdots,\Delta D_K\}$ 为 $D$ 的一个分割。

\noindent\textbf{定义 16.3.1}\quad 设 $S\subset\mathbb R^3$ 是光滑曲面, 函数 $f(x,y,z)$ 在 $S$ 上有定义, 又设\allowbreak $T=\{\Delta S_1,\Delta S_2,\cdots,\Delta S_K\}$ 是 $S$ 的一个分割, 且仍用 $\Delta S_k(k=1,2,\cdots,K)$ 来记每块小曲面的面积, 并记
\[
  \lambda(T)=\max_{1\le k\le K}\operatorname{diam}(\Delta S_k).
\]
在 $\Delta S_k$ 上任取一点 $(\xi_k,\eta_k,\zeta_k)$, 作和式 $\sum\limits_{k=1}^K f(\xi_k,\eta_k,\zeta_k)\Delta S_k$. 若存在常数 $I$, 使得对于 $\forall\varepsilon>0$, $\exists\delta>0$, 对于 $S$ 的任意分割 $T$ 及任取的 $(\xi_k,\eta_k,\zeta_k)\in\Delta S_k(k=1,2,\cdots,K)$, 当 $\lambda(T)<\delta$ 时, 有
\[
  \left|\sum_{k=1}^K f(\xi_k,\eta_k,\zeta_k)\Delta S_k-I\right|<\varepsilon,
\]
即
\[
  \lim_{\lambda(T)\to0}\sum_{k=1}^K f(\xi_k,\eta_k,\zeta_k)\Delta S_k=I,
\]
则称 $f(x,y,z)$ 在 $S$ 上的第一型曲面积分存在, 并称 $I$ 为 $f(x,y,z)$ 在 $S$ 上的第一型曲面积分, 记为
\[
  I=\iint_S f(x,y,z)\mathrm dS,
\]
其中 $f(x,y,z)$ 称为被积函数, $S$ 称为积分曲面。

从定义可以推出第一型曲面积分具有以下性质:

(1) 若在 $S$ 上函数 $f(x,y,z)\equiv1$, 则 $\iint_S f(x,y,z)\mathrm dS=\iint_S\mathrm dS$ 即为曲面 $S$ 的面积。

(2) 若曲面 $S\subset\mathbb R^2$, 则函数 $f(x,y)$ 在 $S$ 上的曲面积分即为二重积分 $\iint_S f(x,y)\mathrm dx\mathrm dy$。

另外, 第一型曲面积分同样具有关于被积函数的线性性以及积分曲面的可加性等性质, 请读者自己给出上述性质的精确描述。

\subsection{第一型曲面积分的存在性与计算公式}

对于光滑曲面上连续函数的第一型曲面积分, 我们有以下的存在性定理以及计算公式。

\noindent\textbf{定理 16.3.2}\quad 设 $S\subset\mathbb R^3$ 是光滑曲面, 其参数方程为
\[
  \begin{cases}
    x=x(u,v),\\
    y=y(u,v),\\
    z=z(u,v),
  \end{cases}
  \quad (u,v)\in D,
\]
其中 $D$ 是由有限条光滑曲线围成的有界闭区域, 又设函数 $f(x,y,z)$ 在 $S$ 上连续, 则 $f(x,y,z)$ 在 $S$ 上的第一型曲面积分存在, 并且
\begin{equation}
  \iint_S f(x,y,z)\mathrm dS
  =\iint_D f(x(u,v),y(u,v),z(u,v))\sqrt{EG-F^2}\mathrm du\mathrm dv,
  \tag{16.3.4}
\end{equation}
其中 $E,G,F$ 由式 (16.3.3) 定义。

\noindent\textbf{证明}\quad 考查曲面 $S$ 的任一分割 $T=\{\Delta S_1,\Delta S_2,\cdots,\Delta S_K\}$, 记其相应于 $D$ 的分割为 $\Delta=\{\Delta D_1,\Delta D_2,\cdots,\Delta D_K\}$. 在每个 $\Delta S_k(k=1,2,\cdots,K)$ 上任取点 $(\xi_k,\zeta_k,\eta_k)$, 则存在 $(u_k,v_k)\in\Delta D_k$, 使得
\[
  (\xi_k,\zeta_k,\eta_k)=(x(u_k,v_k),y(u_k,v_k),z(u_k,v_k)).
\]
因为
\[
  \Delta S_k=\iint_{\Delta D_k}\sqrt{EG-F^2}\mathrm du\mathrm dv,
\]
所以由重积分第一中值定理得
\[
  \Delta S_k=\sqrt{EG-F^2}\bigg|_{(u'_k,v'_k)}\Delta\sigma_k,
\]
其中 $(u'_k,v'_k)\in\Delta D_k,\Delta\sigma_k$ 为 $\Delta D_k$ 的面积. 因此
\[
\begin{aligned}
  &\sum_{k=1}^K f(\xi_k,\zeta_k,\eta_k)\Delta S_k\\
  &\quad=\sum_{k=1}^K f(x(u_k,v_k),y(u_k,v_k),z(u_k,v_k))
    \sqrt{EG-F^2}\bigg|_{(u'_k,v'_k)}\Delta\sigma_k\\
  &\quad=\sum_{k=1}^K f(x(u_k,v_k),y(u_k,v_k),z(u_k,v_k))
    \sqrt{EG-F^2}\bigg|_{(u_k,v_k)}\Delta\sigma_k\\
  &\qquad+\bigg[
    \sum_{k=1}^K f(x(u_k,v_k),y(u_k,v_k),z(u_k,v_k))
    \sqrt{EG-F^2}\bigg|_{(u'_k,v'_k)}\Delta\sigma_k\\
  &\qquad\quad-
    \sum_{k=1}^K f(x(u_k,v_k),y(u_k,v_k),z(u_k,v_k))
    \sqrt{EG-F^2}\bigg|_{(u_k,v_k)}\Delta\sigma_k
  \bigg]\\
  &\quad\triangleq I_1+I_2.
\end{aligned}
\]
显然
\[
  \lim_{\lambda(T)\to0}I_1=\lim_{\lambda(\Delta)\to0}I_1
  =\iint_D f(x(u,v),y(u,v),z(u,v))\sqrt{EG-F^2}\mathrm du\mathrm dv.
\]
记 $M$ 为 $|f(x,y,z)|$ 在 $S$ 上的最大值, $M_k,m_k$ 分别为连续函数 $\sqrt{EG-F^2}$ 在 $\Delta D_k(k=1,2,\cdots,K)$ 上的最大、最小值, 则有
\[
  |I_2|\le M\sum_{k=1}^K(M_k-m_k)\Delta\sigma_k\to0
  \quad(\lambda(\Delta)\to0).
\]
因此
\[
\begin{aligned}
  \lim_{\lambda(T)\to0}\sum_{k=1}^K f(\xi_k,\zeta_k,\eta_k)\Delta S_k
  &=\lim_{\lambda(T)\to0}(I_1+I_2)
    =\lim_{\lambda(\Delta)\to0}(I_1+I_2)\\
  &=\iint_D f(x(u,v),y(u,v),z(u,v))\sqrt{EG-F^2}\mathrm du\mathrm dv.
\end{aligned}
\]
这就证明了 $f(x,y,z)$ 在 $S$ 上的第一型曲面积分存在, 并且
\[
  \iint_S f(x,y,z)\mathrm dS
  =\iint_D f(x(u,v),y(u,v),z(u,v))\sqrt{EG-F^2}\mathrm du\mathrm dv.
\]
证毕。

当曲面 $S$ 由方程 $z=g(x,y)((x,y)\in D)$ 给出时, 由定理 16.3.2 有
\[
  \iint_S f(x,y,z)\mathrm dS
  =\iint_D f(x,y,g(x,y))\sqrt{1+g_x'^2(x,y)+g_y'^2(x,y)}\mathrm dx\mathrm dy.
\]

\noindent\textbf{例 16.3.3}\quad 求第一型曲面积分 $I=\iint_S x^2z\mathrm dS$, 其中 $S$ 为以下的曲面:

(1) 上半单位球面 $z=\sqrt{1-x^2-y^2}$;

(2) 单位球面 $x^2+y^2+z^2=1$。

\noindent\textbf{解}\quad (1) 取 $D=\{(x,y):x^2+y^2\le1\}$, 则
\[
\begin{aligned}
  I&=\iint_S x^2z\mathrm dS\\
  &=\iint_D x^2\cdot\sqrt{1-x^2-y^2}
    \sqrt{1+\frac{x^2}{1-x^2-y^2}+\frac{y^2}{1-x^2-y^2}}\mathrm dx\mathrm dy\\
  &=\iint_D x^2\mathrm dx\mathrm dy
    =\frac12\iint_D(x^2+y^2)\mathrm dx\mathrm dy\\
  &=\frac12\int_0^{2\pi}\mathrm d\theta\int_0^1r^3\mathrm dr=\frac\pi4.
\end{aligned}
\]

(2) 分别记
\[
  S_1:z=\sqrt{1-x^2-y^2},\quad (x,y)\in D,
\]
\[
  S_2:z=-\sqrt{1-x^2-y^2},\quad (x,y)\in D,
\]
则由积分曲面及被积函数的对称性得
\[
  \iint_{S_1}x^2z\mathrm dS=-\iint_{S_2}x^2z\mathrm dS,
\]
所以
\[
  \iint_S x^2z\mathrm dS=0.
\]

\noindent\textbf{例 16.3.4}\quad 求均匀物质曲面 $S:z=f(x,y)=2-(x^2+y^2)(z\ge0)$ 的质心坐标。

\noindent\textbf{解}\quad 设其质心坐标为 $(x_0,y_0,z_0)$, 由对称性有 $x_0=y_0=0$, 并且由微元法有
\[
  z_0=\frac{M_{xy}}{\sigma}=\frac{\iint_S z\mathrm dS}{\iint_S\mathrm dS}
\]
(其中 $M_{xy}$ 是曲面 $S$ 到 $Oxy$ 平面的力矩, $\sigma$ 是曲面 $S$ 的面积). 取 $D=\{(x,y):x^2+y^2\le2\}$. 由 $\sqrt{1+f_x'^2+f_y'^2}=\sqrt{1+4x^2+4y^2}$ 有
\[
\begin{aligned}
  \iint_S\mathrm dS
  &=\iint_D\sqrt{1+4x^2+4y^2}\mathrm dx\mathrm dy
    =\int_0^{2\pi}\mathrm d\theta\int_0^{\sqrt2}\sqrt{1+4r^2}r\mathrm dr\\
  &=2\pi\cdot\frac1{12}(1+4r^2)^{\frac32}\bigg|_0^{\sqrt2}
    =\frac{13}3\pi,
\end{aligned}
\]
\[
\begin{aligned}
  \iint_S z\mathrm dS
  &=\iint_D(2-x^2-y^2)\sqrt{1+4x^2+4y^2}\mathrm dx\mathrm dy\\
  &=\int_0^{2\pi}\mathrm d\theta\int_0^{\sqrt2}r(2-r^2)\sqrt{1+4r^2}\mathrm dr
    =\frac{37}{10}\pi,
\end{aligned}
\]
所以
\[
  z_0=\frac{111}{130}.
\]

\section{第二型曲面积分}

\subsection{曲面的侧}

在日常生活中, 我们常见的曲面大都具有不同的两侧, 如一个自来水管具有内侧与外侧, 当水管破裂时, 管内的水将从内侧流向外侧。

是不是所有的曲面都能分出两侧呢? 著名的麦比乌斯 (Möbius) 带就是一个单侧曲面. 该曲面的构造非常简单: 取一条长方形纸带, 记为 $ABCD$, 将纸带扭转后, 让 $A$ 与 $C$, $B$ 与 $D$ 重合后粘起来, 使纸带构成一个环带 $S$, 该环带则称为麦比乌斯带. 原长方形的边 $AB$ 与 $DC$ 构成该环带 $S$ 的边界. 用颜色来涂该环带, 可以不越过其边界, 而将它全部涂遍颜色. 由此我们无法分出该环带不同的侧 (见图 16.4.1)。

如何用精确的语言来描述一个曲面 $S$ 是单侧还是双侧的呢? 我们有以下定义。

\noindent\textbf{定义 16.4.1}\quad 设 $S$ 是光滑曲面 ($S$ 可以是封闭的, 则此时无边界, 否则它具有边界). 对 $S$ 内部的任一点 $P_0$, 取定 $S$ 在 $P_0$ 处的一个法向的朝向. 若一个动点 $P$ 从 $P_0$ 出发, 不越过 $S$ 的边界, 沿 $S$ 上任何路径 $\Gamma$ 运动回到 $P_0$ 处, $S$ 在 $P$ 的法向从 $P_0$ 处选定的朝向出发连续地沿路径 $\Gamma$ 变化回到 $P_0$ 处时, 总是保持原先在 $P_0$ 处选定的朝向, 则称 $S$ 为双侧曲面, 否则称 $S$ 为单侧曲面。

由定义, 一个单侧曲面 $S$ 上必定存在点 $P_0\in S$, 以及 $S$ 上的一条经过 $P_0$ 的闭路径 $\Gamma$, 若在 $P_0$ 处选择 $S$ 的一个法向的朝向, 然后让 $P$ 从 $P_0$ 出发沿 $\Gamma$ 运动, $P$ 处的法向连续地变化, 当 $P$ 经过 $\Gamma$ 再回到 $P_0$ 时, 法向的朝向与原来确定的朝向相反。

现设 $S$ 是一双侧曲面, 从定义可以看出在 $S$ 上任取一点 $P_0$, 然后取定 $S$ 在 $P_0$ 处一个法向的朝向, 则该曲面一侧法向的朝向即由该处法向的朝向所取定. 这个事实请读者自己来验证. 利用此事实, 我们可以容易地描述由一些参数方程表示的双侧曲面的定侧问题。

例如, 若方程
\[
  \begin{cases}
    x=x(u,v),\\
    y=y(u,v),\\
    z=z(u,v),
  \end{cases}\quad (u,v)\in D
\]
定义了一个光滑的双侧曲面 $S$, 令
\begin{equation}
  A=\frac{\partial(y,z)}{\partial(u,v)},\quad
  B=\frac{\partial(z,x)}{\partial(u,v)},\quad
  C=\frac{\partial(x,y)}{\partial(u,v)},
  \tag{16.4.1}
\end{equation}
则 $\boldsymbol n=\pm(A,B,C)$ 即为 $S$ 的两个不同朝向的法向量. 记 $\boldsymbol n_1=(A,B,C)$, $\boldsymbol n_2=(-A,-B,-C)$. 当 $P_0\in S$, 且在 $P_0$ 处 $C>0$ 时, 我们将由 $\boldsymbol n_1$ 确定的一侧称为 $S$ 的上侧, $\boldsymbol n_2$ 确定的一侧称为 $S$ 下侧. 类似地, 当 $A\ne0$ 时, 我们可以说 $S$ 有前侧和后侧, 而当 $B\ne0$ 时, $S$ 有左侧和右侧。

若光滑曲面 $S$ 的方程为
\[
  z=f(x,y),\quad (x,y)\in D,
\]
则 $S$ 上点 $(x,y,z)$ 处的法向量为 $\boldsymbol n=\pm(-f_x',-f_y',1)$. 因此, 若 $\boldsymbol n$ 取 ``+'' 号, 便得到了 $S$ 的上侧; 若 $\boldsymbol n$ 取 ``-'' 号, 则得到 $S$ 的下侧。

\subsection{第二型曲面积分的定义}

第二型曲面积分的物理背景是流量的计算问题. 设区域 $D\subset\mathbb R^3$ 内有某一流体在流动 (如风的运动, 海水的运动等), 对于 $\forall(x,y,z)\in D$, 其速度函数为
\[
  \boldsymbol v(x,y,z)=(P(x,y,z),Q(x,y,z),R(x,y,z)),
\]
简记为 $\boldsymbol v=(P,Q,R)$. 若不考虑流体的种类, 上述速度函数也可认为是 $D$ 内的一个流速场. 进一步假设 $\boldsymbol v$ 仅仅依赖于 $D$ 内点的位置, 而不依赖于时间的变化, 这样的流速场也称为稳定流速场。

设 $S$ 是 $D$ 内的一个光滑双侧曲面. 我们的问题是: 如何来求单位时间内流体流过 $S$ 的质量? 为了使问题简单化, 我们可设该流体的密度为 1. 因此, 我们只要计算该流体通过 $S$ 的流量即可。

若 $S$ 是平面的一部分, 其面积仍由 $S$ 表示, 而 $\boldsymbol v$ 是常向量函数, 即 $P=C_1,Q=C_2,R=C_3$ 均为常数函数, 记 $S$ 的单位法向量为 $\boldsymbol n$, 则不难看出单位时间内该流体流过 $S$ 的流量
\[
  W=(\boldsymbol v\cdot\boldsymbol n)S,
\]
即流量 $W$ 为 $S$ 的面积与 $\boldsymbol v$ 在 $\boldsymbol n$ 上的投影的乘积. 有了上述的计算公式, 利用积分思想就不难推导出在一般情形下流量的计算公式了。

我们下面仍用微元法来求解此问题. 在 $S$ 上任取一小块 $\Delta S$, 其面积仍然用 $\Delta S$ 记之. 任取 $(\xi,\eta,\zeta)\in\Delta S$, 设 $\Delta S$ 在 $(\xi,\eta,\zeta)$ 处的单位法向量为 $\boldsymbol n(\xi,\eta,\zeta)$, 则可以认为 $\Delta S$ 近似地是法向量为 $\boldsymbol n(\xi,\eta,\zeta)$, 面积为 $\Delta S$ 的一小块平面. 再假设 $\boldsymbol v$ 在 $\Delta S$ 上也近似为 $\boldsymbol v(\xi,\eta,\zeta)$, 则由上面的分析, 我们便得到了流量微元
\[
  \mathrm dW=\boldsymbol v(x,y,z)\cdot\boldsymbol n(x,y,z)\mathrm dS.
\]
因此, 整个流量便是
\[
  W=\iint_S\boldsymbol v\cdot\boldsymbol n\,\mathrm dS
  =\iint_S(P\cos\alpha+Q\cos\beta+R\cos\gamma)\mathrm dS,
\]
其中 $\cos\alpha,\cos\beta,\cos\gamma$ 分别是 $\boldsymbol n$ 的方向余弦. 由于上述积分中含有曲面法向量的方向余弦, 因此它不是第一型曲面积分. 它其实是下面我们要讨论的第二型曲面积分。

对于第二型曲面积分所涉及的曲面 $S$, 我们假定它是分片光滑的双侧曲面. 这类曲面由有限块光滑曲面组成, 任何两块曲面均在边界相交, 并且不存在三块曲面具有相同一段弧作为边界. 如图 16.4.2 中的曲面我们将不予考虑。

现在我们给出分片光滑双侧曲面上的第二型曲面积分的定义。

\noindent\textbf{定义 16.4.2}\quad 设 $S\subset\mathbb R^3$ 是分片光滑双侧曲面, 若它有边界, 则其边界由有限条分段光滑曲线组成. 给定 $S$ 的一侧, $S$ 上每点处的单位法向量记为 $\boldsymbol n=(\cos\alpha,\cos\beta,\cos\gamma)$, 向量函数 $\boldsymbol f(x,y,z)=(P(x,y,z),Q(x,y,z),R(x,y,z))$ 在 $S$ 上有定义. 对 $S$ 作任意分割 $T=\{\Delta S_1,\Delta S_2,\cdots,\Delta S_K\}$, 其中 $\Delta S_k$ 是以光滑曲线为边界的小光滑曲面, 并且仍用 $\Delta S_k(k=1,2,\cdots,K)$ 来记其面积. 记 $\lambda(T)=\max_{1\le k\le K}\operatorname{diam}(\Delta S_k)$. 在 $\Delta S_k$ 上任取一点 $(\xi_k,\eta_k,\zeta_k)(k=1,2,\cdots,K)$. 若存在不依赖于分割 $T$ 以及 $(\xi_k,\eta_k,\zeta_k)(k=1,2,\cdots,K)$ 的选取的常数 $I$, 使得
\begin{equation}
\begin{aligned}
  \lim_{\lambda(T)\to0}\sum_{k=1}^K
  (&P(\xi_k,\eta_k,\zeta_k)\cos\alpha
  +Q(\xi_k,\eta_k,\zeta_k)\cos\beta\\
  &+R(\xi_k,\eta_k,\zeta_k)\cos\gamma)\Delta S_k=I,
\end{aligned}
\tag{16.4.2}
\end{equation}
则称 $I$ 是 $\boldsymbol f=(P,Q,R)$ 在 $S$ 上的第二型曲面积分, 记为
\[
  I=\iint_S P(x,y,z)\mathrm dy\mathrm dz
    +Q(x,y,z)\mathrm dz\mathrm dx
    +R(x,y,z)\mathrm dx\mathrm dy,
\]
其中 $\boldsymbol f=(P,Q,R)$ 称为被积函数, $S$ 称为积分曲面。

\noindent\textbf{注 1}\quad 我们来解释一下第二型曲面积分记号的合理性. 由定义, 一个第二型曲面积分是式 (16.4.2) 定义的极限. 当小曲面 $\Delta S_k$ 的直径很小时, $\Delta S_k$ 近似于 $(\xi_k,\eta_k,\zeta_k)$ 处一小块切平面, 因此它的面积与 $\cos\alpha$ 的乘积近似为该小块切平面在 $Oyz$ 平面上的投影 $\Delta\sigma^k_{yz}$ 的有向面积, 即当 $\cos\alpha>0$ 时, 取其投影的面积为正, 否则为负. 因此, 若仍由 $\Delta\sigma^k_{yz}$ 记其有向面积, 则
\[
\begin{aligned}
  &\lim_{\lambda(T)\to0}\sum_{k=1}^K
    P(\xi_k,\eta_k,\zeta_k)\cos\alpha\Delta S_k\\
  &\quad=\lim_{\lambda(T)\to0}\sum_{k=1}^K
    P(\xi_k,\eta_k,\zeta_k)\Delta\sigma^k_{yz}.
\end{aligned}
\]
上式右边可以自然地记成 $\iint_S P\,\mathrm dy\mathrm dz$. 同理我们可解释记号 $\iint_S Q\,\mathrm dz\mathrm dx$ 和 $\iint_S R\,\mathrm dx\mathrm dy$ 的合理性。

对于第二型曲面积分, 读者不难看出它们仍然具有关于被积函数的线性性以及积分曲面的相加性. 值得注意的是, 当曲面相加时要指定曲面的侧. 特别需要指出的是: 双侧曲面上同一个向量函数在两个侧上的第二型曲面积分的绝对值相等, 但相差一个负号。

\subsection{第二型曲面积分的存在性与计算公式}

设光滑曲面 $S$ 由参数方程
\begin{equation}
  \begin{cases}
    x=x(u,v),\\
    y=y(u,v),\\
    z=z(u,v),
  \end{cases}\quad (u,v)\in D
  \tag{16.4.3}
\end{equation}
给出, 则 $S$ 上每一点处的方向余弦为
\[
  \cos\alpha=\pm\frac{A}{\sqrt{A^2+B^2+C^2}},\quad
  \cos\beta=\pm\frac{B}{\sqrt{A^2+B^2+C^2}},
\]
\[
  \cos\gamma=\pm\frac{C}{\sqrt{A^2+B^2+C^2}},
\]
其中 $A,B,C$ 由式 (16.4.1) 定义. 值得指出的是, 它们都是 $(u,v)\in D$ 的连续函数。

\noindent\textbf{定理 16.4.1}\quad 设 $S$ 为光滑双侧曲面, 其参数方程由 (16.4.3) 给出. 选定 $S$ 的一侧, 记其单位法向量为
\[
  \boldsymbol n=\left(
    \frac{A}{\sqrt{A^2+B^2+C^2}},
    \frac{B}{\sqrt{A^2+B^2+C^2}},
    \frac{C}{\sqrt{A^2+B^2+C^2}}
  \right),
\]
再设向量函数 $\boldsymbol f(x,y,z)=(P,Q,R)$ 在 $S$ 上连续, 则 $\boldsymbol f$ 在 $S$ 上的第二型曲面积分存在, 并且
\[
  \iint_S P\,\mathrm dy\mathrm dz+Q\,\mathrm dz\mathrm dx+R\,\mathrm dx\mathrm dy
  =\iint_D(PA+QB+RC)\,\mathrm du\mathrm dv.
\]

\noindent\textbf{证明}\quad 由 $S$ 的光滑性及式 (16.4.1) 知 $\boldsymbol n$ 在 $S$ 上是连续向量函数. 由
\begin{equation}
  \iint_S P\,\mathrm dy\mathrm dz+Q\,\mathrm dz\mathrm dx+R\,\mathrm dx\mathrm dy
  =\iint_S(P\cos\alpha+Q\cos\beta+R\cos\gamma)\mathrm dS,
  \tag{16.4.4}
\end{equation}
并注意到上式的右边是连续函数 $P\cos\alpha+Q\cos\beta+R\cos\gamma$ 的第一型曲面积分, 因此左边的第二型曲面积分存在. 定理的可积性部分得证。

由公式 (16.3.4) 得
\[
  \iint_S(P\cos\alpha+Q\cos\beta+R\cos\gamma)\mathrm dS
  =\iint_D(PA+QB+RC)\,\mathrm du\mathrm dv.
\]
证毕。

特别地, 若光滑曲面 $S$ 的方程为
\[
  z=f(x,y),\quad (x,y)\in D,
\]
并取定其上侧, 其中 $D\subset\mathbb R^2$ 是由分段光滑曲线所围成的有界闭区域. 由于 $S$ 上的单位法向量平行于 $(-f_x',-f_y',1)$, 容易推出
\[
  \iint_S P(x,y,z)\mathrm dy\mathrm dz+Q(x,y,z)\mathrm dz\mathrm dx+R(x,y,z)\mathrm dx\mathrm dy
\]
\[
  =\iint_D[-P(x,y,f(x,y))f_x'(x,y)-Q(x,y,f(x,y))f_y'(x,y)
  +R(x,y,f(x,y))]\mathrm dx\mathrm dy.
\]

\noindent\textbf{例 16.4.2}\quad 计算第二型曲面积分
\[
  I=\iint_S x\mathrm dy\mathrm dz+y\mathrm dz\mathrm dx+z\mathrm dx\mathrm dy,
\]
其中 $S$ 为由三个坐标平面及平面 $x+y+z=1$ 所围四面体的外侧。

\noindent\textbf{解法 1}\quad 记 $S$ 落在 $Oxy,Oyz,Ozx$ 平面上的部分分别记为 $S_z,S_x$ 及 $S_y$, 落在平面 $x+y+z=1$ 的部分记为 $S_1$ (见图 16.4.3). 在 $S_z$ 上, $z=0$, $\mathrm dy\mathrm dz=\mathrm dz\mathrm dx=0$, 从而
\[
  \iint_{S_z}x\mathrm dy\mathrm dz+y\mathrm dz\mathrm dx+z\mathrm dx\mathrm dy=0.
\]
同理, 在 $S_y$ 与 $S_x$ 上的积分都为零. 因此
\[
  I=\iint_{S_1}x\mathrm dy\mathrm dz+y\mathrm dz\mathrm dx+z\mathrm dx\mathrm dy.
\]
记 $D=\{(x,y):x\ge0,y\ge0,x+y\le1\}$, 则由对称性有
\[
  I=3\iint_D(1-x-y)\mathrm dx\mathrm dy
  =3\int_0^1\mathrm dx\int_0^{1-x}(1-x-y)\mathrm dy=\frac12.
\]

\noindent\textbf{解法 2}\quad 与解法 1 同理得到
\[
  I=\iint_{S_1}x\mathrm dy\mathrm dz+y\mathrm dz\mathrm dx+z\mathrm dx\mathrm dy.
\]
由于在 $S_1$ 上点 $(x,y,z)$ 处的方向余弦 $(\cos\alpha,\cos\beta,\cos\gamma)$ 即为 $S_1$ 的单位法向量, 且
\[
  \cos\alpha=\cos\beta=\cos\gamma=\frac1{\sqrt3},
\]
而 $S_1$ 是一个三角形, 其面积为 $\frac{\sqrt3}{2}$, 因此
\[
  I=\iint_{S_1}\frac{x+y+z}{\sqrt3}\mathrm dS
  =\frac1{\sqrt3}\iint_{S_1}\mathrm dS=\frac12.
\]

\section{各类积分之间的联系}

设 $f(x)\in C^1[a,b]$. 众所周知, 牛顿-莱布尼茨公式
\[
  \int_a^b\mathrm df(x)=f(b)-f(a)
\]
是一元微积分中最重要的公式. 换种观点来看该公式, 如果我们定义一个函数 $f(x)$ 在 $a,b(a<b)$ 两点上的积分为 $f(b)-f(a)$, 则上述公式建立了区间 $[a,b]$ 上 $f'(x)$ 的积分与 $f(x)$ 在端点处的积分的联系. 因此, 一个自然的问题是: 是否平面或空间中一个几何体的边界与内部上的某些积分总存在联系? 本节中我们将讨论这个问题。

\subsection{格林公式}

格林 (Green) 公式给出的是平面区域上函数的二重积分与其相关的函数在边界上的第二型曲线积分之间的联系. 在这里, 我们只讨论 $\mathbb R^2$ 中的有界闭区域 $D$, 其边界是有限条约当曲线, 且每条约当曲线都是分段光滑的. 因此 $D$ 一定具有如图 16.5.1 的形状。

对于平面 $\mathbb R^2$ 内一条约当曲线 $\Gamma$, 以 $\Gamma$ 为边界的区域有两个: 一个是有界区域, 该区域也称为 $\Gamma$ 的内部; 而另一个是无界区域, 该区域也称为 $\Gamma$ 的外部. 如在图 16.5.1 中, $D$ 位于 $\Gamma_0$ 的内部, 而在 $\Gamma_k(k=1,2,\cdots,K)$ 的外部. 一个区域 $D$ 称为单连通区域, 如果位于 $D$ 内的任一条约当曲线的内部总包含在 $D$ 内; 否则称 $D$ 是多连通区域. 在我们假定的情形下, 若 $D$ 是一条约当曲线的内部, 则它是单连通的; 若 $D$ 由 $K(K>1)$ 条约当曲线围成, 则 $D$ 是多连通的, 此时也称 $D$ 是 $K$ 连通的。

一条约当曲线具有两个互为相反的方向. 现设 $D$ 是由有限条约当曲线所围成的区域. 我们在第二册曾给出 $\partial D$ 的定向, 即若一个人沿着 $D$ 的边界曲线上的一个方向前进时, 曲线所围的区域总在他的左边, 则称该方向为 $\partial D$ 的正向. 如图 16.5.1 所示, 相对 $D$ 来说, $\Gamma_0$ 的正向是逆时针方向, 而 $\Gamma_k(k=1,2,\cdots,K)$ 的正向是顺时针方向。

有了上述准备, 我们可以证明下述定理。

\noindent\textbf{定理 16.5.1(格林公式)}\quad 设 $D\subset\mathbb R^2$ 是有界闭区域, 其边界由有限条分段光滑的约当曲线组成. 若函数 $P(x,y),Q(x,y)$ 在 $D$ 上具有连续偏导数, 则有
\[
  \int_{\partial D}P\mathrm dx+Q\mathrm dy
  =\iint_D\left(\frac{\partial Q}{\partial x}-\frac{\partial P}{\partial y}\right)\mathrm dx\mathrm dy,
\]
其中 $\partial D$ 的方向为正向。

\noindent\textbf{证明}\quad 我们先证 $D$ 是单连通区域的情形. 证明分以下几步进行:

(1) 设 $D$ 是 $Y$ 型区域, 即
\[
  D=\{(x,y):\varphi(y)\le x\le\psi(y),c\le y\le d\},
\]
如图 16.5.2 所示. 下面我们证明
\[
  \int_{\partial D}Q(x,y)\mathrm dy=\iint_D\frac{\partial Q}{\partial x}\mathrm dx\mathrm dy.
\]
我们有
\[
  \iint_D\frac{\partial Q}{\partial x}\mathrm dx\mathrm dy
  =\int_c^d\mathrm dy\int_{\varphi(y)}^{\psi(y)}
  \frac{\partial Q}{\partial x}\mathrm dx
  =\int_c^d[Q(\psi(y),y)-Q(\varphi(y),y)]\mathrm dy.
\]
由
\[
\begin{aligned}
  \int_{\partial D}Q(x,y)\mathrm dy
  &=\int_{\widehat{AB}}Q(x,y)\mathrm dy+\int_{\widehat{BC}}Q(x,y)\mathrm dy
    +\int_{\widehat{CE}}Q(x,y)\mathrm dy+\int_{\widehat{EA}}Q(x,y)\mathrm dy,\\
  \int_{\widehat{AB}}Q(x,y)\mathrm dy&=\int_{\widehat{CE}}Q(x,y)\mathrm dy=0,\\
  \int_{\widehat{BC}}Q(x,y)\mathrm dy&=\int_c^d Q(\psi(y),y)\mathrm dy,
\end{aligned}
\]
及
\[
  \int_{\widehat{EA}}Q(x,y)\mathrm dy
  =\int_d^cQ(\varphi(y),y)\mathrm dy
  =-\int_c^dQ(\varphi(y),y)\mathrm dy,
\]
因此有
\[
  \int_{\partial D}Q(x,y)\mathrm dy=\iint_D\frac{\partial Q}{\partial x}\mathrm dx\mathrm dy.
\]

(2) 若 $D$ 是 $X$ 型区域, 与 (1) 同理可证
\[
  \int_{\partial D}P(x,y)\mathrm dx=-\iint_D\frac{\partial P}{\partial y}\mathrm dx\mathrm dy.
\]

(3) 若 $D$ 既是 $X$ 型又是 $Y$ 型区域, 则成立
\[
  \int_{\partial D}P\mathrm dx+Q\mathrm dy
  =\iint_D\left(\frac{\partial Q}{\partial x}-\frac{\partial P}{\partial y}\right)\mathrm dx\mathrm dy.
\]

若区域 $D$ 加上有限条光滑曲线后, 可分成有限个既是 $X$ 型又是 $Y$ 型区域, 则可在每个小区域上应用上述结果. 将每个小区域边界上的曲线积分与区域上的重积分分别相加, 由于添加的曲线必然是两个小区域的公共边界, 因此在曲线积分求和中出现两次, 但因方向相反而相互抵消, 所以总有
\[
  \int_{\partial D}P\mathrm dx+Q\mathrm dy
  =\iint_D\left(\frac{\partial Q}{\partial x}-\frac{\partial P}{\partial y}\right)\mathrm dx\mathrm dy.
\]
对一般的单连通区域, 可以用分成有限个既是 $X$ 型又是 $Y$ 型区域的区域逼近方法证之, 但这方面内容在数学分析课程中难以介绍, 因此我们略去该证明。

当 $D$ 是 $K(K>1)$ 连通时, 我们同样可以在区域 $D$ 内添加若干条光滑曲线, 将 $D$ 分成有限个单连通区域 (见图 16.5.3). 在每个小单连通区域上由上所证成立格林公式, 然后将每个小区域边界上的曲线积分与区域上的重积分分别相加, 注意到所添加的光滑曲线仍然是两个小区域的公共边界, 因此和式中在所添加曲线上的曲线积分将相互抵消, 所以格林公式对 $K$ 连通区域仍然成立. 证毕。

在上述格林公式中, 我们假定了函数 $P$ 和 $Q$ 在 $D$ 上具有连续偏导数, 这就要求 $P,Q$ 在包含 $D$ 的某个开集内具有连续偏导数. 另外, 对于有限条封闭曲线 $\Gamma$ 所围的区域 $D$, 为了强调曲线的封闭性, 今后我们常常用
\[
  \oint_{\partial D}P\mathrm dx+Q\mathrm dy
\]
来表示沿 $\partial D$ 关于 $D$ 的正向的积分。

对于格林公式, 我们还可以有以下形式. 设光滑曲线 $\Gamma$ 由
\[
  \begin{cases}
    x=x(t),\\
    y=y(t),
  \end{cases}\quad (t\in[\alpha,\beta])
\]
给出, $t$ 的增加方向恰好是 $\Gamma$ 的正向. 记该曲线在 $(x,y)$ 处的单位切向量为 $\boldsymbol v$, 则有
\[
  \mathrm dx=x'(t)\mathrm dt
  =\frac{x'(t)}{\sqrt{x'^2(t)+y'^2(t)}}\sqrt{x'^2(t)+y'^2(t)}\mathrm dt
  =\cos(\boldsymbol v,x)\mathrm ds,
\]
其中 $\cos(\boldsymbol v,x)$ 为切向量 $\boldsymbol v$ 与 $x$ 轴正向的夹角的余弦. 同理, 有
\[
  \mathrm dy=y'(t)\mathrm dt=\cos(\boldsymbol v,y)\mathrm ds,
\]
其中 $\cos(\boldsymbol v,y)$ 为切向量 $\boldsymbol v$ 与 $y$ 轴正向的夹角的余弦. 设函数 $P(x,y),Q(x,y)$ 在 $\Gamma$ 上连续, 则
\[
  \int_\Gamma P\mathrm dx+Q\mathrm dy
  =\int_\Gamma[P\cos(\boldsymbol v,x)+Q\cos(\boldsymbol v,y)]\mathrm ds.
\]
现取 $\Gamma$ 在 $(x,y)$ 处的单位法向量为 $\boldsymbol n$, 使得 $\boldsymbol n$ 与 $\boldsymbol v$ 成右手系. 这样, 在封闭曲线 $\Gamma$ 围成区域 $D$ 时, 若取 $\Gamma$ 的正向, 则选取的法向量指向 $D$ 的外部 (见图 16.5.4), 从而有
\[
  \begin{cases}
    \cos(\boldsymbol n,x)=\cos(\boldsymbol v,y),\\
    \cos(\boldsymbol n,y)=-\cos(\boldsymbol v,x),
  \end{cases}
\]
其中 $\cos(\boldsymbol n,x)$ 与 $\cos(\boldsymbol n,y)$ 分别为法向 $\boldsymbol n$ 与 $x$ 轴, $y$ 轴正向的夹角的余弦. 因此
\[
\begin{aligned}
  &\int_\Gamma[P\cos(\boldsymbol n,x)+Q\cos(\boldsymbol n,y)]\mathrm ds\\
  &=\int_\Gamma[-Q\cos(\boldsymbol v,x)+P\cos(\boldsymbol v,y)]\mathrm ds\\
  &=\int_\Gamma(-Q)\mathrm dx+P\mathrm dy
  =\iint_D\left(\frac{\partial P}{\partial x}+\frac{\partial Q}{\partial y}\right)\mathrm dx\mathrm dy.
\end{aligned}
\]

另外, 在定积分的应用中, 我们曾用
\[
  \int_\Gamma x\mathrm dy=-\int_\Gamma y\mathrm dx
  =\frac12\int_\Gamma x\mathrm dy-y\mathrm dx
\]
来表示封闭曲线 $\Gamma$ 所围区域 $D$ 的面积. 现在利用格林公式可以清楚看到上述三个积分都等于 $\iint_D\mathrm dx\mathrm dy$, 即 $D$ 的面积。

\noindent\textbf{例 16.5.1}\quad 设 $\Gamma=\widehat{AB}\subset\mathbb R^2$ 是从点 $A(0,0)$ 到点 $B(0,1)$ 的任一条光滑曲线, 试计算第二型曲线积分
\[
  \int_\Gamma 2xy\mathrm dx+(x^2-y^2)\mathrm dy.
\]

\noindent\textbf{解}\quad 由于 $\Gamma$ 是任一条光滑曲线, 因此直接计算将无法求出积分值. 注意到若令
\[
  P(x,y)=2xy,\quad Q(x,y)=x^2-y^2,
\]
则有
\[
  \frac{\partial Q}{\partial x}=2x=\frac{\partial P}{\partial y}.
\]
因此对于 $\forall(x,y)\in\mathbb R^2$, 有 $\frac{\partial Q}{\partial x}-\frac{\partial P}{\partial y}=0$。

为了应用格林公式, 设实轴上从点 $(0,0)$ 到点 $(0,1)$ 的线段为 $\Gamma_2$. 我们再取一条光滑曲线 $\Gamma_1$, 使得 $\Gamma_1$ 和 $\Gamma_2$ 都以点 $(0,0),(0,1)$ 为端点, 并且 $\Gamma\cup\Gamma_1^-$ 和 $\Gamma_1\cup\Gamma_2^-$ 均为约当区域 (见图 16.5.5). 设 $\Gamma\cup\Gamma_1^-$ 所围的有界闭区域为 $D_1$, 而 $\Gamma_1\cup\Gamma_2^-$ 所围的有界闭区域为 $D_2$, 由格林公式有
\[
\begin{aligned}
  \int_\Gamma P\mathrm dx+Q\mathrm dy
  &=\int_{\Gamma_1}P\mathrm dx+Q\mathrm dy+\int_{\Gamma\cup\Gamma_1^-}P\mathrm dx+Q\mathrm dy\\
  &=\int_{\Gamma_1}P\mathrm dx+Q\mathrm dy+\iint_{D_1}0\mathrm dx\mathrm dy\\
  &=\int_{\Gamma_2}P\mathrm dx+Q\mathrm dy+\int_{\Gamma_1\cup\Gamma_2^-}P\mathrm dx+Q\mathrm dy\\
  &=\int_{\Gamma_2}P\mathrm dx+Q\mathrm dy+\iint_{D_2}0\mathrm dx\mathrm dy\\
  &=\int_0^1 0\mathrm dx=0.
\end{aligned}
\]

\noindent\textbf{例 16.5.2}\quad 计算第二型曲线积分
\[
  \oint_\Gamma\frac{x\mathrm dy-y\mathrm dx}{x^2+y^2},
\]
其中 $\Gamma\subset\mathbb R^2$ 为任一条光滑约当曲线, 且 $(0,0)\notin\Gamma$。

\noindent\textbf{解}\quad 记
\[
  P(x,y)=\frac{-y}{x^2+y^2},\quad Q(x,y)=\frac{x}{x^2+y^2},
\]
则对任意的 $(x,y)\ne(0,0)$, 有
\[
  \frac{\partial Q}{\partial x}=\frac{y^2-x^2}{(x^2+y^2)^2}=\frac{\partial P}{\partial y}.
\]
因此, 设 $\Gamma$ 所围成的区域为 $D$, 则当 $(0,0)\notin D$ 时, 由格林公式有
\[
  \oint_\Gamma\frac{x\mathrm dy-y\mathrm dx}{x^2+y^2}
  =\iint_D0\mathrm dx\mathrm dy=0.
\]
当 $(0,0)\in D$ 时, 格林公式条件不满足. 对于 $R>0$, 设
\[
  \Gamma_R=\{(x,y):x^2+y^2=R^2\}.
\]
取 $R$ 充分大, 使得 $\Gamma$ 落在 $\Gamma_R$ 的内部. 如图 16.5.6, 则 $\Gamma$ 与 $\Gamma_R$ 围成一个二连通区域 $D_R$, 在此区域上应用格林公式, 有
\[
  0=\iint_{D_R}\left(\frac{\partial Q}{\partial x}-\frac{\partial P}{\partial y}\right)\mathrm dx\mathrm dy
  =\int_{\Gamma_R}P\mathrm dx+Q\mathrm dy+\int_{\Gamma^-}P\mathrm dx+Q\mathrm dy.
\]
因此, 我们有
\[
\begin{aligned}
  \int_\Gamma P\mathrm dx+Q\mathrm dy
  &=\int_{\Gamma_R}\frac{x\mathrm dy-y\mathrm dx}{x^2+y^2}
    =\frac1{R^2}\int_{\Gamma_R}x\mathrm dy-y\mathrm dx\\
  &=\frac2{R^2}\iint_{D_R}\mathrm dx\mathrm dy=2\pi.
\end{aligned}
\]

\subsection{高斯公式}

在本小节中, 设 $D\subset\mathbb R^3$ 是有界闭区域, $\partial D$ 由有限块光滑双侧曲面所组成. 选定 $\partial D$ 的侧为 $D$ 的外侧, 即 $D$ 的外法向量确定的侧 (此时也称外侧为 $\partial D$ 的正向). 高斯 (Gauss) 公式将给出函数在 $D$ 上的三重积分与相关函数在 $\partial D$ 上的曲面积分的关系。

\noindent\textbf{定理 16.5.2(高斯公式)}\quad 设 $D\subset\mathbb R^3$ 是有界闭区域, $\partial D$ 是 $D$ 的外侧, 函数\allowbreak $P(x,y,z),Q(x,y,z),R(x,y,z)$ 在 $D$ 上具有连续偏导数, 则有
\[
  \iint_{\partial D}P\mathrm dy\mathrm dz+Q\mathrm dz\mathrm dx+R\mathrm dx\mathrm dy
  =\iiint_D\left(\frac{\partial P}{\partial x}+\frac{\partial Q}{\partial y}
  +\frac{\partial R}{\partial z}\right)\mathrm dx\mathrm dy\mathrm dz.
\]

\noindent\textbf{证明}\quad 高斯公式的证明与格林公式的证明具有相似性, 基本思想是: 先设 $D$ 是一些特殊的区域, 然后再用区域逼近来证明一般区域上的高斯公式. 同样, 区域逼近这一步在这里也不作介绍。

(1) 设 $D$ 是如下的闭区域:
\[
  D=\{(x,y,z):(y,z)\in\Omega,\varphi(y,z)\le x\le\psi(y,z)\}
\]
(通常也称为 $X$ 型区域), 其中 $\Omega\subset\mathbb R^2$ 是由有限多条光滑闭曲线所围的有界闭区域, $\varphi(y,z)$ 与 $\psi(y,z)$ 是 $\Omega$ 上的连续函数. 我们来证明
\[
  \iint_{\partial D}P(x,y,z)\mathrm dy\mathrm dz
  =\iiint_D\frac{\partial P}{\partial x}\mathrm dx\mathrm dy\mathrm dz.
\]
其证明方法是将上式两边的积分都化为 $\Omega$ 上的二重积分. 记
\[
\begin{aligned}
  S_1&=\{(x,y,z):x=\psi(y,z),(y,z)\in\Omega\},\quad \text{取前侧};\\
  S_2&=\{(x,y,z):x=\varphi(y,z),(y,z)\in\Omega\},\quad \text{取后侧};\\
  S_3&=\{(x,y,z):\varphi(y,z)\le x\le\psi(y,z),(y,z)\in\Omega\},
    \quad \text{取 $S_3$ 的侧为 $S$ 的外侧在 $S_3$ 的限制}.
\end{aligned}
\]
由于
\[
  \iint_{S_1}P(x,y,z)\mathrm dy\mathrm dz
  =\iint_\Omega P(\psi(y,z),y,z)\mathrm dy\mathrm dz,
\]
\[
  \iint_{S_2}P(x,y,z)\mathrm dy\mathrm dz
  =-\iint_\Omega P(\varphi(y,z),y,z)\mathrm dy\mathrm dz,\quad
  \iint_{S_3}P(x,y,z)\mathrm dy\mathrm dz=0,
\]
因此
\[
  \iint_{\partial D}P(x,y,z)\mathrm dy\mathrm dz
  =\iint_\Omega[P(\psi(y,z),y,z)-P(\varphi(y,z),y,z)]\mathrm dy\mathrm dz.
\]
而
\[
\begin{aligned}
  \iiint_D\frac{\partial P}{\partial x}\mathrm dx\mathrm dy\mathrm dz
  &=\iint_\Omega\mathrm dy\mathrm dz
    \int_{\varphi(y,z)}^{\psi(y,z)}\frac{\partial P}{\partial x}\mathrm dx\\
  &=\iint_\Omega[P(\psi(y,z),y,z)-P(\varphi(y,z),y,z)]\mathrm dy\mathrm dz,
\end{aligned}
\]
所以
\[
  \iint_{\partial D}P(x,y,z)\mathrm dy\mathrm dz
  =\iiint_D\frac{\partial P}{\partial x}\mathrm dx\mathrm dy\mathrm dz.
\]

(2) 若 $D=\{(x,y,z):(x,y)\in\Omega,\varphi(x,y)\le z\le\psi(x,y)\}$ (称之为 $Z$ 型区域, 其中 $\Omega,\varphi,\psi$ 所满足的条件与 (1) 中的相同), 则我们可以证明
\[
  \iint_{\partial D}R(x,y,z)\mathrm dx\mathrm dy
  =\iiint_D\frac{\partial R}{\partial z}\mathrm dx\mathrm dy\mathrm dz.
\]

(3) 当 $D=\{(x,y,z):(z,x)\in\Omega,\varphi(z,x)\le y\le\psi(z,x)\}$ (也称为 $Y$ 型区域, 其中 $\Omega,\varphi,\psi$ 所满足的条件与 (1) 中的相同) 时, 我们可以证明
\[
  \iint_{\partial D}Q(x,y,z)\mathrm dz\mathrm dx
  =\iiint_D\frac{\partial Q}{\partial y}\mathrm dx\mathrm dy\mathrm dz.
\]

(4) 显然, 当 $D$ 刚好同时是这三种区域时, 定理结论成立. 对于一般的中间无 ``洞'' 的有界闭区域 $D$, 当区域 $D$ 可以用有限块光滑曲面将其分成一些小闭区域, 且使得每个小闭区域同时是 $X$ 型, $Y$ 型与 $Z$ 型区域时 (如四面体区域), 高斯公式在每个小区域上成立. 再将每个小区域上相应积分相加, 注意到所添加曲面必是两个小区域的公共边界, 因此在第二型曲面积分中出现两次, 但因它们的侧刚好相反, 从而相互抵消. 所以对这类区域, 高斯公式依然成立. 若 $D$ 不能分成有限块同时是上述的 $X$ 型, $Y$ 型和 $Z$ 型区域, 我们不加证明地指出, 利用分成同时是上述三种区域的区域逼近方法, 可以证明对于 $D$ 仍成立高斯公式。

(5) 若 $D$ 是中间有 ``洞'' 的情形. 此时, 我们亦可添加若干块分片光滑的曲面, 将 $D$ 分成有限个中间无 ``洞'' 的小区域. 由上所证, 高斯公式在这些小区域上成立. 注意到所添加曲面必是两个小区域的公共边界, 因此对这类相当广泛的区域, 高斯公式依然成立. 证毕。

\noindent\textbf{注}\quad 对于区域有 ``洞'' 的情形, 区域 $D$ 内的 ``洞'' $D_1$ 的边界曲面的侧相对于区域 $D$ 来说, 应该是 $\partial D_1$ 的内侧 (见图 16.5.7)。

\noindent\textbf{例 16.5.3}\quad 计算第二型曲面积分
\[
  I=\iint_S(\sin yz+x)\mathrm dy\mathrm dz+(e^{xz}+y)\mathrm dz\mathrm dx+(xy+z)\mathrm dx\mathrm dy,
\]
其中 $S$ 是单位球面上半部分的上侧, 即
\[
  S=\{(x,y,z):x^2+y^2+z^2=1,z\ge0\}.
\]

\noindent\textbf{解}\quad 直接计算似乎比较复杂, 因此利用高斯公式来计算. 由于 $S$ 不封闭, 若我们取 $Oxy$ 平面上的单位圆盘 $S_1=\{(x,y):x^2+y^2\le1\}$, 并取 $S_1$ 的下侧, 则 $S\cup S_1$ 构成了上半单位球体 $D$ 的边界, 且取外侧. 由高斯公式得
\[
\begin{aligned}
  &\iint_{S\cup S_1}(\sin yz+x)\mathrm dy\mathrm dz+(e^{xz}+y)\mathrm dz\mathrm dx+(xy+z)\mathrm dx\mathrm dy\\
  &=3\iiint_D\mathrm dx\mathrm dy\mathrm dz
    =3\cdot\frac12\cdot\frac{4\pi}{3}=2\pi,
\end{aligned}
\]
而在 $S_1$ 上 $\mathrm dy\mathrm dz=0,\mathrm dz\mathrm dx=0$, 又
\[
  \iint_{S_1}(xy+z)\mathrm dx\mathrm dy
  =-\iint_{x^2+y^2\le1}xy\mathrm dx\mathrm dy=0,
\]
所以 $I=2\pi$。

\noindent\textbf{例 16.5.4}\quad 设 $S\subset\mathbb R^3$ 为封闭光滑曲面, 取外侧, 以它为边界的有界闭区域是 $D$. 已知点 $(\xi,\eta,\zeta)\in\mathbb R^3$ 不在 $S$ 上. 计算高斯积分
\[
  I=\iint_S\frac{\cos(\boldsymbol r,\boldsymbol n)}{r^2}\mathrm dS,
\]
其中 $\boldsymbol r=(x-\xi,y-\eta,z-\zeta)$, $r=|\boldsymbol r|$, $\boldsymbol n$ 是 $S$ 的单位外法向量。

\noindent\textbf{解}\quad 由于
\[
  \cos(\boldsymbol r,\boldsymbol n)
  =\frac{x-\xi}{r}\cos(\boldsymbol n,x)
   +\frac{y-\eta}{r}\cos(\boldsymbol n,y)
   +\frac{z-\zeta}{r}\cos(\boldsymbol n,z),
\]
因此
\[
  I=\iint_S\frac{x-\xi}{r^3}\mathrm dy\mathrm dz
    +\frac{y-\eta}{r^3}\mathrm dz\mathrm dx
    +\frac{z-\zeta}{r^3}\mathrm dx\mathrm dy.
\]
因为
\[
  \frac{\partial}{\partial x}\left(\frac{x-\xi}{r^3}\right)=\frac1{r^3}-\frac{3(x-\xi)^2}{r^5},\quad
  \frac{\partial}{\partial y}\left(\frac{y-\eta}{r^3}\right)=\frac1{r^3}-\frac{3(y-\eta)^2}{r^5},
\]
\[
  \frac{\partial}{\partial z}\left(\frac{z-\zeta}{r^3}\right)=\frac1{r^3}-\frac{3(z-\zeta)^2}{r^5},
\]
所以
\[
  \frac{\partial}{\partial x}\left(\frac{x-\xi}{r^3}\right)
  +\frac{\partial}{\partial y}\left(\frac{y-\eta}{r^3}\right)
  +\frac{\partial}{\partial z}\left(\frac{z-\zeta}{r^3}\right)=0.
\]
因此, 当 $(\xi,\eta,\zeta)\notin D$ 时, 由高斯公式得 $I=0$。

当 $(\xi,\eta,\zeta)\in D$ 时, 我们可取 $\varepsilon$ 充分小, 使得球面
\[
  S_\varepsilon=\{(x,y,z):(x-\xi)^2+(y-\eta)^2+(z-\zeta)^2=\varepsilon^2\}
\]
完全落在 $D$ 的内部. 取 $S_\varepsilon$ 的内侧 $S_\varepsilon^-$, 设闭区域 $D_\varepsilon$ 以 $S\cup S_\varepsilon^-$ 为边界, 则
\[
  \iint_{S\cup S_\varepsilon^-}\frac{\cos(\boldsymbol r,\boldsymbol n)}{r^2}\mathrm dS
  =\iiint_{D_\varepsilon}0\mathrm dx\mathrm dy\mathrm dz=0.
\]
注意到在 $S_\varepsilon$ 上, $\boldsymbol r$ 与 $\boldsymbol n$ 平行, 所以
\[
\begin{aligned}
  \iint_S\frac{\cos(\boldsymbol r,\boldsymbol n)}{r^2}\mathrm dS
  &=-\iint_{S_\varepsilon^-}\frac{\cos(\boldsymbol r,\boldsymbol n)}{r^2}\mathrm dS
    =\iint_{S_\varepsilon}\frac{\cos(\boldsymbol r,\boldsymbol n)}{r^2}\mathrm dS\\
  &=\iint_{S_\varepsilon}\frac{\mathrm dS}{\varepsilon^2}
    =\frac1{\varepsilon^2}4\pi\varepsilon^2=4\pi.
\end{aligned}
\]

\subsection{斯托克斯公式}

斯托克斯 (Stokes) 公式是格林公式的直接推广, 它揭示了函数在空间曲面 $S$ 上的第二型曲面积分与相关函数在 $\partial S$ 上的第二型曲线积分之间的关系. 为此我们先讨论一下空间曲面与其边界的定向问题。

设 $S$ 是空间分片光滑的双侧曲面, 其边界由有限条分段光滑的空间闭曲线组成. 给定 $S$ 的一侧, 设其由法向量 $\boldsymbol n$ 所确定. 若一个人站立与该法向量保持一致并且沿 $\partial S$ 前进时 $S$ 在其左边, 则称他的前进方向为 $\partial S$ 关于该曲面确定的侧的正向。

如图 16.5.8, $\Gamma_0$ 的逆时针方向及 $\Gamma_k(k=1,2,\cdots,K)$ 的顺时针方向构成了 $S$ 关于 $\boldsymbol n$ 确定的侧的正向. 显然, 这样的定向与平面区域的边界定向是保持一致的。

\noindent\textbf{定理 16.5.3(斯托克斯公式)}\quad 设 $S\subset\mathbb R^3$ 是光滑双侧曲面, $\partial S$ 由有限多条分段光滑曲线组成, 给定 $S$ 的一侧并取 $\partial S$ 关于该侧为正向的方向, 再设函数 $P(x,y,z),Q(x,y,z),R(x,y,z)$ 在包含 $S$ 的某个区域内具有连续偏导数, 则
\[
\begin{aligned}
  \int_{\partial S}P\mathrm dx+Q\mathrm dy+R\mathrm dz
  &=\iint_S\left(\frac{\partial R}{\partial y}-\frac{\partial Q}{\partial z}\right)\mathrm dy\mathrm dz\\
  &\quad+\left(\frac{\partial P}{\partial z}-\frac{\partial R}{\partial x}\right)\mathrm dz\mathrm dx
    +\left(\frac{\partial Q}{\partial x}-\frac{\partial P}{\partial y}\right)\mathrm dx\mathrm dy.
\end{aligned}
\]

\noindent\textbf{证明}\quad 若 $\partial S$ 由有限条闭曲线所组成, 我们可以在 $S$ 上添加若干条曲线将 $S$ 分成有限块小曲面, 而每块小曲面的边界只有一条闭曲线. 如同格林公式的证明, 若在小曲面上斯托克斯公式成立, 则 $S$ 上也成立. 因此不妨设 $S$ 的边界是一条分段光滑的闭曲线. 另外, 我们也只对几种简单的曲面进行证明, 对于一般的曲面, 可以由简单曲面逼近来证明, 但这个过程非常复杂, 在此就不作介绍了。

具体地说, 我们先设 $S$ 由方程
\begin{equation}
  z=f(x,y),\quad (x,y)\in D
  \tag{16.5.1}
\end{equation}
给出, 其中 $D\subset\mathbb R^2$ 是由一条分段光滑的约当曲线所围成的区域, $\partial D$ 的定向由 $S$ 在 $D$ 的投影所确定。

一方面, 由第二型曲线积分的计算公式以及格林公式得
\[
\begin{aligned}
  \int_{\partial S}P(x,y,z)\mathrm dx
  &=\int_{\partial D}P(x,y,f(x,y))\mathrm dx\\
  &=-\iint_D\left(\frac{\partial P}{\partial y}
    +\frac{\partial P}{\partial z}\cdot\frac{\partial f}{\partial y}\right)\mathrm dx\mathrm dy.
\end{aligned}
\]
另一方面, 有
\[
\begin{aligned}
  \iint_S\frac{\partial P}{\partial z}\mathrm dz\mathrm dx
  -\frac{\partial P}{\partial y}\mathrm dx\mathrm dy
  &=\iint_D\left[\frac{\partial P}{\partial z}\left(-\frac{\partial f}{\partial y}\right)
    -\frac{\partial P}{\partial y}\right]\mathrm dx\mathrm dy\\
  &=-\iint_D\left(\frac{\partial P}{\partial y}
    +\frac{\partial P}{\partial z}\cdot\frac{\partial f}{\partial y}\right)\mathrm dx\mathrm dy.
\end{aligned}
\]
因此
\[
  \int_{\partial S}P\mathrm dx
  =\iint_S\frac{\partial P}{\partial z}\mathrm dz\mathrm dx
    -\frac{\partial P}{\partial y}\mathrm dx\mathrm dy.
\]
同样的计算可以得出
\[
  \int_{\partial S}Q\mathrm dy
  =\iint_S\left(-\frac{\partial Q}{\partial z}\right)\mathrm dy\mathrm dz
    +\frac{\partial Q}{\partial x}\mathrm dx\mathrm dy.
\]

同理, 对于曲面 $S$ 由方程 $x=g(y,z)((y,z)\in D)$ 给出的情形, 其中 $S$ 与 $D$ 的光滑性假设如 (16.5.1), 则有
\[
  \int_{\partial S}Q\mathrm dy+\int_{\partial S}R\mathrm dz
  =\iint_S\left(\frac{\partial R}{\partial y}-\frac{\partial Q}{\partial z}\right)\mathrm dy\mathrm dz
    +\frac{\partial Q}{\partial x}\mathrm dx\mathrm dy-\frac{\partial R}{\partial x}\mathrm dz\mathrm dx;
\]
而对于曲面 $S$ 由方程 $y=h(z,x)((z,x)\in D)$ 给出的情形, 则有
\[
  \int_{\partial S}P\mathrm dx+\int_{\partial S}R\mathrm dz
  =\iint_S\left(\frac{\partial P}{\partial z}-\frac{\partial R}{\partial x}\right)\mathrm dz\mathrm dx
    +\frac{\partial R}{\partial y}\mathrm dy\mathrm dz-\frac{\partial P}{\partial y}\mathrm dx\mathrm dy.
\]
因此, 若 $S$ 是同时能表示成上述三种形式的曲面或用光滑曲线可以将 $S$ 分成有限块这样的曲面, 则有
\[
\begin{aligned}
  \int_{\partial S}P\mathrm dx+Q\mathrm dy+R\mathrm dz
  &=\iint_S\left(\frac{\partial R}{\partial y}-\frac{\partial Q}{\partial z}\right)\mathrm dy\mathrm dz\\
  &\quad+\left(\frac{\partial P}{\partial z}-\frac{\partial R}{\partial x}\right)\mathrm dz\mathrm dx
    +\left(\frac{\partial Q}{\partial x}-\frac{\partial P}{\partial y}\right)\mathrm dx\mathrm dy.
\end{aligned}
\]
证毕。

在斯托克斯公式中, 曲面积分的被积函数似乎不太容易记忆, 我们可以用下述方式来记该公式:
\[
  \int_{\partial S}P\mathrm dx+Q\mathrm dy+R\mathrm dz
  =\iint_S
  \begin{vmatrix}
    \cos\alpha&\cos\beta&\cos\gamma\\
    \frac{\partial}{\partial x}&\frac{\partial}{\partial y}&\frac{\partial}{\partial z}\\
    P&Q&R
  \end{vmatrix}\mathrm dS
\]
\[
  =\iint_S
  \begin{vmatrix}
    \mathrm dy\mathrm dz&\mathrm dz\mathrm dx&\mathrm dx\mathrm dy\\
    \frac{\partial}{\partial x}&\frac{\partial}{\partial y}&\frac{\partial}{\partial z}\\
    P&Q&R
  \end{vmatrix},
\]
其中 $(\cos\alpha,\cos\beta,\cos\gamma)$ 为与 $\partial S$ 的正向保持一致的法向量的方向余弦. 另外, 由于此时 $\partial S$ 为封闭曲线, 我们也用 $\oint_{\partial S}P\mathrm dx+Q\mathrm dy+R\mathrm dz$ 来记 $\int_{\partial S}P\mathrm dx+Q\mathrm dy+R\mathrm dz$。

\noindent\textbf{例 16.5.5}\quad 计算第二型曲线积分
\[
  I=\oint_{\Gamma_h}(y^2-z^2)\mathrm dx+(z^2-x^2)\mathrm dy+(x^2-y^2)\mathrm dz,
\]
其中 $\Gamma_h$ 是平面 $x+y+z=h(-1<h<1)$ 与球面 $x^2+y^2+z^2=1$ 的交线, 从 $z$ 轴看去取逆时针方向。

\noindent\textbf{解}\quad 设平面 $x+y+z=h(-1<h<1)$ 被圆周 $\Gamma_h$ 所围的部分为 $S_h$, 则 $S_h$ 是半径为 $\sqrt{1-h^2}$ 的圆盘. 取 $S_h$ 的法向量向上, 即选取法向量的方向余弦为
\[
  (\cos\alpha,\cos\beta,\cos\gamma)=\left(\frac1{\sqrt3},\frac1{\sqrt3},\frac1{\sqrt3}\right).
\]
利用斯托克斯公式, 有
\[
  I=\iint_{S_h}
  \begin{vmatrix}
    \frac1{\sqrt3}&\frac1{\sqrt3}&\frac1{\sqrt3}\\
    \frac{\partial}{\partial x}&\frac{\partial}{\partial y}&\frac{\partial}{\partial z}\\
    y^2-z^2&z^2-x^2&x^2-y^2
  \end{vmatrix}\mathrm dS
\]
\[
  =-\frac4{\sqrt3}\iint_{S_h}(x+y+z)\mathrm dS
  =-\frac4{\sqrt3}\iint_{S_h}h\mathrm dS
  =-\frac{4\pi}{\sqrt3}h(1-h^2).
\]

\section{微分形式简介}

\subsection{微分形式}

我们前面学过了许多不同的积分, 对一个多元函数来说, 笼统地说它的积分是十分不确切的事情; 就三元向量函数而言, 给了一个三元向量函数 $(P(x,y,z),Q(x,y,z),R(x,y,z))$, 对它就可以进行曲线积分与曲面积分. 但如果我们写出 $P\mathrm dx+Q\mathrm dy+R\mathrm dz$ 或 $P\mathrm dy\mathrm dz+Q\mathrm dz\mathrm dx+R\mathrm dx\mathrm dy$, 就可以清楚知道前者可以在曲线上积分, 而后者可以在曲面上积分. 而 $P\mathrm dx+Q\mathrm dy+R\mathrm dz$ 与 $P\mathrm dy\mathrm dz+Q\mathrm dz\mathrm dx+R\mathrm dx\mathrm dy$ 就是所谓的微分形式. 当然, 微分形式的意义远不止这些, 它是现代数学中一种基本的工具, 在许多领域中有着广泛的应用. 由于在后续课程中有专门课程来介绍它, 在这里我们只作一些简单介绍, 主要目的是使读者能对重积分换元公式及各种积分之间的联系有更深刻的理解。

设 $f(\boldsymbol x)=f(x_1,x_2,\cdots,x_n)$ 是定义在 $\mathbb R^n(n\ge2)$ 中的一个可微函数, 则它的微分是
\[
  \mathrm df=\sum_{i=1}^n\frac{\partial f(\boldsymbol x)}{\partial x_i}\mathrm dx_i.
\]
现在抛开微分的几何意义, 纯粹从数学的观点来看上式右边的表达式, 我们可以认为 $\sum_{i=1}^n\frac{\partial f(\boldsymbol x)}{\partial x_i}\mathrm dx_i$ 是 $\mathbb R^n$ 中点 $\boldsymbol x$ 处的一个关于 $(\mathrm dx_1,\mathrm dx_2,\cdots,\mathrm dx_n)$ 的线性组合. 因此, 我们可以将 $(\mathrm dx_1,\mathrm dx_2,\cdots,\mathrm dx_n)$ 看做一组基. 对给定 $\mathbb R^n$ 中的一个区域 $D$, $C^1(D)$ ($D$ 内连续可微函数的全体) 在这组基上的线性组合
\[
  \sum_{i=1}^n a_i(\boldsymbol x)\mathrm dx_i\quad
  (a_i(\boldsymbol x)\in C^1(D),\ i=1,2,\cdots,n)
\]
称为 $D$ 内的一个一次微分形式或 $1$-形式, 一般用 $\omega,\eta$ 等来记之. 所有 $D$ 内的 $1$-形式的全体记为 $\Lambda^1(D)$. 在后面的讨论中, 我们总是固定区域 $D$, 因此 $\Lambda^1(D)$ 也记为 $\Lambda^1$。

在 $\Lambda^1$ 内, 我们可以自然地定义加法与数乘如下:

(1) 设 $\omega_j=\sum_{i=1}^n a_i^j(\boldsymbol x)\mathrm dx_i\in\Lambda^1(j=1,2)$, 定义
\[
  \omega_1+\omega_2=\sum_{i=1}^n(a_i^1(\boldsymbol x)+a_i^2(\boldsymbol x))\mathrm dx_i\in\Lambda^1.
\]

(2) 设 $\omega=\sum_{i=1}^n a_i(\boldsymbol x)\mathrm dx_i\in\Lambda^1$, $\lambda\in\mathbb R$, 定义
\[
  \lambda\omega=\sum_{i=1}^n(\lambda a_i(\boldsymbol x))\mathrm dx_i\in\Lambda^1,
\]
另外, 我们规定 $0=\sum_{i=1}^n0\mathrm dx_i$。

读者不难验证上述两种运算满足一个线性空间所要求的各种运算定律, 从而 $\Lambda^1$ 构成一个线性空间。

现对任意的 $k(1<k\le n)$, 在 $(\mathrm dx_1,\mathrm dx_2,\cdots,\mathrm dx_n)$ 中取 $k$ 个元素, 将其形式地记成 $\mathrm dx_{i_1}\wedge\mathrm dx_{i_2}\wedge\cdots\wedge\mathrm dx_{i_k}$, 这里 $\wedge$ 暂且作为一个形式记号. 然后规定
\[
  \mathrm dx_{i_1}\wedge\cdots\wedge\mathrm dx_{i_j}\wedge\mathrm dx_{i_{j+1}}\wedge\cdots\wedge\mathrm dx_{i_k}
  =-\mathrm dx_{i_1}\wedge\cdots\wedge\mathrm dx_{i_{j+1}}\wedge\mathrm dx_{i_j}\wedge\cdots\wedge\mathrm dx_{i_k}.
\]
在此规定下, 若 $i_j=i_l(1\le j<l\le k)$, 则 $\mathrm dx_{i_1}\wedge\mathrm dx_{i_2}\wedge\cdots\wedge\mathrm dx_{i_k}=0$. 因此我们有 $C_n^k$ 个有序元
\[
  \mathrm dx_{i_1}\wedge\mathrm dx_{i_2}\wedge\cdots\wedge\mathrm dx_{i_k}
  \quad (1\le i_1<i_2<\cdots<i_k\le n).
\]
我们将以这 $C_n^k$ 个有序元作为基并且系数在 $C^1(D)$ 的线性空间记为 $\Lambda^k$. $\Lambda^k$ 中的元素称为 $k$ 次微分形式, 一般我们还是用 $\omega,\eta$ 等来记它们. 从定义我们知道, 对于 $\forall\omega\in\Lambda^k$, 我们有下述表达式:
\[
  \omega=\sum_{1\le i_1<i_2<\cdots<i_k\le n}
  a_{i_1i_2\cdots i_k}(\boldsymbol x)
  \mathrm dx_{i_1}\wedge\mathrm dx_{i_2}\wedge\cdots\wedge\mathrm dx_{i_k},
\]
其中 $a_{i_1i_2\cdots i_k}(\boldsymbol x)\in C^1(D)$. $\omega$ 的这种表示称为它的标准形式. 特别地, 我们将 $C^1(D)$ 记为 $\Lambda^0$, 即 $\Lambda^0$ 是 $D$ 内连续可微函数的全体, 它的基是 $f(\boldsymbol x)=1$, 从而它是一维线性空间. 另外, 我们记 $0$ 为任意次零微分形式。

当 $k=n$ 时, $\Lambda^n$ 只有一个基 $\mathrm dx_1\wedge\mathrm dx_2\wedge\cdots\wedge\mathrm dx_n$, 从而 $\Lambda^n$ 也是一维线性空间。

在 $\mathbb R^2$ 中可记 $(x_1,x_2)=(x,y)$. 在 $\mathbb R^3$ 中, 我们记 $(x_1,x_2,x_3)=(x,y,z)$. 例如, 在 $\mathbb R^3$ 中设 $\omega\in\Lambda^2$, 则 $\omega$ 的标准形式为
\[
  \omega=P\mathrm dy\wedge\mathrm dz+Q\mathrm dx\wedge\mathrm dz+R\mathrm dx\wedge\mathrm dy,
\]
其中 $P,Q,R$ 是 $\mathbb R^3$ 中的连续可微函数。

\subsection{微分形式的外积}

外积运算建立了不同次微分形式空间的联系. 一个 $p$ 次微分形式 $\omega$ 与一个 $q$ 次微分形式 $\eta$ 的外积记为 $\omega\wedge\eta$, 它是一个 $p+q$ 次微分形式. 因此只有当 $p+q\le n$ 时, 外积才有意义. 外积 $\omega\wedge\eta$ 的确切定义为: 设
\[
  \omega=\sum_{1\le i_1<i_2<\cdots<i_p\le n}
  a_{i_1i_2\cdots i_p}(\boldsymbol x)
  \mathrm dx_{i_1}\wedge\mathrm dx_{i_2}\wedge\cdots\wedge\mathrm dx_{i_p}\in\Lambda^p,
\]
\[
  \eta=\sum_{1\le j_1<j_2<\cdots<j_q\le n}
  a_{j_1j_2\cdots j_q}(\boldsymbol x)
  \mathrm dx_{j_1}\wedge\mathrm dx_{j_2}\wedge\cdots\wedge\mathrm dx_{j_q}\in\Lambda^q,
\]
则
\[
  \omega\wedge\eta=
  \sum_{\substack{1\le i_1<i_2<\cdots<i_p\le n\\
  1\le j_1<j_2<\cdots<j_q\le n}}
  a_{i_1i_2\cdots i_p}(\boldsymbol x)a_{j_1j_2\cdots j_q}(\boldsymbol x)
  \cdot\mathrm dx_{i_1}\wedge\mathrm dx_{i_2}\wedge\cdots\wedge\mathrm dx_{i_p}
  \wedge\mathrm dx_{j_1}\wedge\mathrm dx_{j_2}\wedge\cdots\wedge\mathrm dx_{j_q}\in\Lambda^{p+q}.
\]
特别地, 当 $f\in\Lambda^0,\omega\in\Lambda^k$ 时, $f\wedge\omega$ 定义为 $f\omega$。

容易验证, 外积具有下述性质:

(1) 设 $f_1,f_2\in\Lambda^0,\omega_1,\omega_2\in\Lambda^p,\eta\in\Lambda^q$, 则有
\[
  (f_1\omega_1+f_2\omega_2)\wedge\eta
  =f_1\omega_1\wedge\eta+f_2\omega_2\wedge\eta.
\]

(2) 设 $\omega\in\Lambda^p,\eta\in\Lambda^q$, 则有 $\omega\wedge\eta=(-1)^{pq}\eta\wedge\omega$。

上述性质 (2) 也称为外积的反对称性. 此性质的证明可由一个微分形式中的基元素两两换位改变符号直接推出。

(3) 设 $\omega,\eta,\theta$ 是任意三个微分形式, 则 $\omega\wedge(\eta\wedge\theta)=(\omega\wedge\eta)\wedge\theta$。

有了外积的概念, 微分形式 $\mathrm dx_1\wedge\mathrm dx_2$ 可以看成是 $\mathrm dx_1$ 与 $\mathrm dx_2$ 的外积, 类似地理解 $\mathrm dx_1\wedge\mathrm dx_2\wedge\cdots\wedge\mathrm dx_n$。

\noindent\textbf{例 16.6.1}\quad 设
\[
  \begin{pmatrix}x\\y\end{pmatrix}
  =
  \begin{pmatrix}x(u,v)\\y(u,v)\end{pmatrix}
\]
是 $D\subset\mathbb R^2$ 到 $\Omega\subset\mathbb R^2$ 的 $C^1$ 同胚映射, 求 $\mathrm dx\wedge\mathrm dy$ 关于 $\mathrm du\wedge\mathrm dv$ 的表示。

\noindent\textbf{解}\quad 由微分定义有
\[
  \mathrm dx=\frac{\partial x}{\partial u}\mathrm du+\frac{\partial x}{\partial v}\mathrm dv,
  \quad
  \mathrm dy=\frac{\partial y}{\partial u}\mathrm du+\frac{\partial y}{\partial v}\mathrm dv,
\]
因此
\[
  \mathrm dx\wedge\mathrm dy
  =\left(\frac{\partial x}{\partial u}\mathrm du+\frac{\partial x}{\partial v}\mathrm dv\right)
  \wedge
  \left(\frac{\partial y}{\partial u}\mathrm du+\frac{\partial y}{\partial v}\mathrm dv\right)
\]
\[
  =\left(\frac{\partial x}{\partial u}\frac{\partial y}{\partial v}
  -\frac{\partial x}{\partial v}\frac{\partial y}{\partial u}\right)\mathrm du\wedge\mathrm dv
  =\frac{\partial(x,y)}{\partial(u,v)}\mathrm du\wedge\mathrm dv.
\]

\noindent\textbf{例 16.6.2}\quad 设
\[
  \begin{pmatrix}x\\y\\z\end{pmatrix}
  =
  \begin{pmatrix}x(u,v,w)\\y(u,v,w)\\z(u,v,w)\end{pmatrix}
\]
是 $D\subset\mathbb R^3$ 到 $\Omega\subset\mathbb R^3$ 的 $C^1$ 同胚映射, 求 $\mathrm dx\wedge\mathrm dy\wedge\mathrm dz$ 关于 $\mathrm du\wedge\mathrm dv\wedge\mathrm dw$ 的表示。

\noindent\textbf{解}\quad 因为
\[
  \mathrm dx=\frac{\partial x}{\partial u}\mathrm du+\frac{\partial x}{\partial v}\mathrm dv+\frac{\partial x}{\partial w}\mathrm dw,
\]
\[
  \mathrm dy=\frac{\partial y}{\partial u}\mathrm du+\frac{\partial y}{\partial v}\mathrm dv+\frac{\partial y}{\partial w}\mathrm dw,
\]
\[
  \mathrm dz=\frac{\partial z}{\partial u}\mathrm du+\frac{\partial z}{\partial v}\mathrm dv+\frac{\partial z}{\partial w}\mathrm dw,
\]
所以
\[
  \mathrm dx\wedge\mathrm dy
  =\left(\frac{\partial x}{\partial u}\mathrm du+\frac{\partial x}{\partial v}\mathrm dv+\frac{\partial x}{\partial w}\mathrm dw\right)
  \wedge
  \left(\frac{\partial y}{\partial u}\mathrm du+\frac{\partial y}{\partial v}\mathrm dv+\frac{\partial y}{\partial w}\mathrm dw\right)
\]
\[
  =\frac{\partial(x,y)}{\partial(v,w)}\mathrm dv\wedge\mathrm dw
  +\frac{\partial(x,y)}{\partial(w,u)}\mathrm dw\wedge\mathrm du
  +\frac{\partial(x,y)}{\partial(u,v)}\mathrm du\wedge\mathrm dv.
\]
\[
  \mathrm dx\wedge\mathrm dy\wedge\mathrm dz
  =\frac{\partial z}{\partial u}\cdot\frac{\partial(x,y)}{\partial(v,w)}
  \mathrm dv\wedge\mathrm dw\wedge\mathrm du
  +\frac{\partial z}{\partial v}\cdot\frac{\partial(x,y)}{\partial(w,u)}
  \mathrm dw\wedge\mathrm du\wedge\mathrm dv
\]
\[
  \quad+\frac{\partial z}{\partial w}\cdot\frac{\partial(x,y)}{\partial(u,v)}
  \mathrm du\wedge\mathrm dv\wedge\mathrm dw
  =\frac{\partial(x,y,z)}{\partial(u,v,w)}\mathrm du\wedge\mathrm dv\wedge\mathrm dw.
\]

下面我们以区域 $D\subset\mathbb R^2$ 到 $\Omega\subset\mathbb R^2$ 的 $C^1$ 同胚映射
\[
  \begin{pmatrix}x\\y\end{pmatrix}
  =
  \begin{pmatrix}x(u,v)\\y(u,v)\end{pmatrix}
\]
来说明其雅可比行列式的符号与映射的保向性的联系. 记 $\boldsymbol u=(u,v),\boldsymbol x=(x,y)$, 并将上述同胚映射记成 $\boldsymbol x=\boldsymbol f(\boldsymbol u)$。

任取一点 $\boldsymbol u_0\in D$, 则由 $C^1$ 同胚映射的性质, 存在 $U(\boldsymbol u_0,\delta_0)\subset D$, 使得 $\boldsymbol f(U(\boldsymbol u_0,\delta_0))$ 是包含 $\boldsymbol f(\boldsymbol u_0)=\boldsymbol x_0$ 的一个区域. 因此任取一条约当曲线 $\Gamma\subset U(\boldsymbol u_0,\delta_0)$, 则 $\gamma=\boldsymbol f(\Gamma)$ 是 $\Omega$ 内的一条约当曲线. 如果 $\boldsymbol x=\boldsymbol f(\boldsymbol u)$ 将 $\Gamma$ 的正向映成 $\gamma$ 的正向 (即动点 $\boldsymbol u$ 沿 $\Gamma$ 逆时针运动时, 动点 $\boldsymbol f(\boldsymbol u)$ 沿 $\gamma$ 逆时针运动), 则称 $\boldsymbol f(\boldsymbol u)$ 在 $\boldsymbol u_0$ 处保向. 若 $\boldsymbol f(\boldsymbol u)$ 在 $\boldsymbol u_0$ 处将 $\Gamma$ 的正向映成 $\gamma$ 的负向, 则称 $\boldsymbol f(\boldsymbol u)$ 在 $\boldsymbol u_0$ 处反定向. 可以证明, 若 $\boldsymbol f(\boldsymbol u)$ 在 $\boldsymbol u_0$ 处保向, 则它在任意点 $\boldsymbol u\in D$ 都是保向的 (见本章习题)。

下面我们来研究保向映射的雅可比行列式的性质. 任取 $\boldsymbol u_0=(u_0,v_0)\in D$, 任取过 $\boldsymbol u_0$ 的一条光滑曲线
\[
  \Gamma:\begin{cases}
    u=u(t),\\
    v=v(t),
  \end{cases}\quad t\in(\alpha,\beta),
\]
使得它满足 $u(t_0)=u_0,v(t_0)=v_0(t_0\in(\alpha,\beta))$. $\Gamma$ 在 $\boldsymbol u_0$ 处的切线斜率为
\[
  k_u=\frac{v'(t_0)}{u'(t_0)}.
\]
曲线 $\gamma=\boldsymbol f(\Gamma)$ 的参数方程为
\[
  \begin{cases}
    x=x(u(t),v(t)),\\
    y=y(u(t),v(t)).
  \end{cases}
\]
它在 $\boldsymbol x_0=\boldsymbol f(\boldsymbol u_0)$ 处的切线的斜率
\[
  k_x=\frac{y'(t_0)}{x'(t_0)}
  =\frac{y_u'(\boldsymbol x_0)u'(t_0)+y_v'(\boldsymbol x_0)v'(t_0)}
  {x_u'(\boldsymbol x_0)u'(t_0)+x_v'(\boldsymbol x_0)v'(t_0)}
  =\frac{y_u'(\boldsymbol x_0)+y_v'(\boldsymbol x_0)k_u}
  {x_u'(\boldsymbol x_0)+x_v'(\boldsymbol x_0)k_u}.
\]
所以
\[
  \frac{\mathrm dk_x}{\mathrm dk_u}
  =
  \frac{\left.\dfrac{\partial(x,y)}{\partial(u,v)}\right|_{\boldsymbol u_0}}
  {[x_u'(\boldsymbol x_0)+x_v'(\boldsymbol x_0)k_u]^2}.
\]
从上式可以看出, 当 $\left.\dfrac{\partial(x,y)}{\partial(u,v)}\right|_{\boldsymbol u_0}>0$ 时, $k_x$ 是 $k_u$ 的单调上升函数. 由导数的几何意义知, 此时 $\boldsymbol f(\boldsymbol u)$ 在 $\boldsymbol u_0$ 处保向; 当 $\left.\dfrac{\partial(x,y)}{\partial(u,v)}\right|_{\boldsymbol u_0}<0$ 时, $\boldsymbol f(\boldsymbol u)$ 在 $\boldsymbol u_0$ 处反定向. 我们已经知道, 一个变换在一点保向则处处保向, 因此 $\dfrac{\partial(x,y)}{\partial(u,v)}$ 在 $D$ 内必定不变号。

现在回到微分形式的外积. 设
\[
  \begin{pmatrix}x\\y\end{pmatrix}
  =
  \begin{pmatrix}x(u,v)\\y(u,v)\end{pmatrix}
\]
是区域 $D\subset\mathbb R^2$ 到可求面积的有界闭区域 $\Omega\subset\mathbb R^2$ 的 $C^1$ 同胚映射. 当它保向时, 有 $\dfrac{\partial(x,y)}{\partial(u,v)}\ge0$. 这时我们有
\[
  \mathrm dx\wedge\mathrm dy
  =\frac{\partial(x,y)}{\partial(u,v)}\mathrm du\wedge\mathrm dv.
\]
当变换
\[
  \begin{pmatrix}x\\y\end{pmatrix}
  =
  \begin{pmatrix}x(u,v)\\y(u,v)\end{pmatrix}
\]
不保向时, 则 $\dfrac{\partial(x,y)}{\partial(u,v)}\le0$. 现在我们把 $Ouv$ 平面上的 $u$ 轴改为 $v$ 轴, 而把 $v$ 轴改成 $u$ 轴. 即令 $u'=v,v'=u$, 则上述变换雅可比行列式为
\[
  \frac{\partial(x,y)}{\partial(u',v')}
  =
  \begin{vmatrix}
    \frac{\partial x}{\partial u'}&\frac{\partial x}{\partial v'}\\
    \frac{\partial y}{\partial u'}&\frac{\partial y}{\partial v'}
  \end{vmatrix}
  =
  \begin{vmatrix}
    \frac{\partial x}{\partial v}&\frac{\partial x}{\partial u}\\
    \frac{\partial y}{\partial v}&\frac{\partial y}{\partial u}
  \end{vmatrix}
  =-\frac{\partial(x,y)}{\partial(u,v)}\ge0,
\]
从而
\[
  \mathrm dx\wedge\mathrm dy
  =\frac{\partial(x,y)}{\partial(u',v')}\mathrm du'\wedge\mathrm dv'
  =\frac{\partial(x,y)}{\partial(u,v)}\mathrm du\wedge\mathrm dv.
\]
因此, 若 $f(x,y)$ 是 $\Omega$ 上的可积函数, 当映射
\[
  \begin{pmatrix}x\\y\end{pmatrix}
  =
  \begin{pmatrix}x(u,v)\\y(u,v)\end{pmatrix}
\]
保向时, 有
\[
  \iint_\Omega f(x,y)\mathrm dx\wedge\mathrm dy
  =
  \iint_D f(x(u,v),y(u,v))\frac{\partial(x,y)}{\partial(u,v)}\mathrm du\wedge\mathrm dv.
\]
而当
\[
  \begin{pmatrix}x\\y\end{pmatrix}
  =
  \begin{pmatrix}x(u,v)\\y(u,v)\end{pmatrix}
\]
不保向时, 仍然成立
\[
  \iint_\Omega f(x,y)\mathrm dx\wedge\mathrm dy
  =
  \iint_D f(x(u,v),y(u,v))\frac{\partial(x,y)}{\partial(u,v)}\mathrm du\wedge\mathrm dv.
\]
同理, 对 $n$ 重积分也有类似的结果. 这说明了引入微分形式的外积后, 我们在重积分的变量替换中不必考虑变换的雅可比行列式的正负号. 当 $n\ge3$ 时, 一个映射的保向性的刻画比 $\mathbb R^2$ 的情形会更加复杂, 在这里我们就不作介绍了。

\subsection{外微分}

微分形式的外微分是一个重要的运算, 利用它可以将许多我们学过的重要定理统一起来. 下面我们先对 $\mathbb R^n(n\ge2)$ 中区域 $D$ 内的微分形式引进外微分这一概念。

\noindent\textbf{定义 16.6.1}\quad 设 $\omega\in\Lambda^k(1\le k\le n)$, 其形式为
\[
  \omega=
  \sum_{1\le i_1<i_2<\cdots<i_k\le n}
  a_{i_1i_2\cdots i_k}(\boldsymbol x)
  \mathrm dx_{i_1}\wedge\mathrm dx_{i_2}\wedge\cdots\wedge\mathrm dx_{i_k},
\]
称 $k+1$ 次微分形式
\[
  \sum_{1\le i_1<i_2<\cdots<i_k\le n}
  \left(\sum_{j=1}^n\frac{\partial a_{i_1i_2\cdots i_k}(\boldsymbol x)}{\partial x_j}\mathrm dx_j\right)
  \wedge\mathrm dx_{i_1}\wedge\mathrm dx_{i_2}\wedge\cdots\wedge\mathrm dx_{i_k}
\]
为 $\omega$ 的外微分, 记为 $\mathrm d\omega$。

显然, 当 $\omega\in\Lambda^n$ 时, $\mathrm d\omega=0$. 当 $\omega\in\Lambda^k$, 且 $\omega$ 在 $\Lambda^k$ 的基元素前面的系数函数是 $D$ 内的 $C^1$ 函数时, 称 $\omega$ 是 $C^1$ 微分形式, 这时 $\mathrm d\omega$ 的系数仅在 $D$ 内连续. 若要讨论 $\mathrm d(\mathrm d\omega)\triangleq\mathrm d^2\omega$ (称之为二次外微分, 即外微分的外微分), 则要假定 $\omega$ 是 $C^2$ 微分形式, 即 $\omega$ 中在基前面的系数函数在 $D$ 内具有各个二阶连续偏导数。

由定义可推知外微分具有下述性质:

(1) 设 $\omega_1,\omega_2\in\Lambda^k$, 则 $\mathrm d(\omega_1+\omega_2)=\mathrm d\omega_1+\mathrm d\omega_2$。

(2) 设 $\omega\in\Lambda^k,\eta\in\Lambda^l$, 则 $\mathrm d(\omega\wedge\eta)=\mathrm d\omega\wedge\eta+(-1)^k\omega\wedge\mathrm d\eta$。

(3) 当 $\omega\in\Lambda^k$, 且 $\omega$ 是 $C^2$ 微分形式时, $\mathrm d(\mathrm d\omega)=\mathrm d^2\omega=0$。

这些性质的证明留给读者。

下面我们来求一些微分形式的外微分。

(1) 设 $\omega=P(x,y)\mathrm dx+Q(x,y)\mathrm dy$, 则
\[
  \mathrm d\omega=\left(\frac{\partial Q}{\partial x}-\frac{\partial P}{\partial y}\right)\mathrm dx\wedge\mathrm dy.
\]

(2) 设 $\eta=P(x,y,z)\mathrm dx+Q(x,y,z)\mathrm dy+R(x,y,z)\mathrm dz$, 则
\[
  \mathrm d\eta=
  \left(\frac{\partial R}{\partial y}-\frac{\partial Q}{\partial z}\right)\mathrm dy\wedge\mathrm dz
  +
  \left(\frac{\partial P}{\partial z}-\frac{\partial R}{\partial x}\right)\mathrm dz\wedge\mathrm dx
  +
  \left(\frac{\partial Q}{\partial x}-\frac{\partial P}{\partial y}\right)\mathrm dx\wedge\mathrm dy.
\]

(3) 设 $\theta=P(x,y,z)\mathrm dy\wedge\mathrm dz+Q(x,y,z)\mathrm dz\wedge\mathrm dx+R(x,y,z)\mathrm dx\wedge\mathrm dy$, 则
\[
  \mathrm d\theta=
  \left(\frac{\partial P}{\partial x}+\frac{\partial Q}{\partial y}+\frac{\partial R}{\partial z}\right)
  \mathrm dx\wedge\mathrm dy\wedge\mathrm dz.
\]

本节开始时我们曾指出, 给了 $\mathbb R^3$ 中的一个微分形式, 我们很清楚它在 $\mathbb R^3$ 中什么样的子集上可积分. 因此, 对于一个几何体 $D$ (闭区域、曲面、曲线等), 由上述外微分的表达式, 格林公式、高斯公式以及斯托克斯公式可以统一地写成
\[
  \int_D\mathrm d\omega=\int_{\partial D}\omega,
\]
其中 $\omega$ 是 $\overline D$ 上的一个 $C^1$ 微分形式, $\partial D$ 取正向。

若在牛顿-莱布尼茨公式 $\int_a^b f'(x)\mathrm dx=f(x)\big|_a^b$ 中将 $f(x)\big|_a^b=f(b)-f(a)$ 看成是 $\omega=f(x)$ 在 $a,b$ 两点的积分, 上述公式也可写成 $\int_D\mathrm d\omega=\int_{\partial D}\omega$ 的形式, 其中 $D=[a,b],\partial D=\{a,b\}$。

公式 $\int_D\mathrm d\omega=\int_{\partial D}\omega$ 也称为斯托克斯公式. 尽管如此, 任何积分公式最后的计算以及它们的证明都必须要用到牛顿-莱布尼茨公式, 所以人们还是称牛顿-莱布尼茨公式为微积分基本定理。

另外, 当 $n\ge4$ 时, 微分形式具有更多的形式, 因此在 $\mathbb R^n$ 中具有更多可以讨论积分问题的子集. 关于这些问题在微分流形课程中将做系统的介绍, 在这里我们就不展开这方面的内容了。

\section{曲线积分与路径的无关性}

设 $D\subset\mathbb R^3$ 是一个区域, $u(x,y,z)$ 是 $D$ 内的 $C^1$ 函数, 则 $\mathrm du=u_x'\mathrm dx+u_y'\mathrm dy+u_z'\mathrm dz$ 是 $u(x,y,z)$ 在 $D$ 内的全微分. 给定 $A_0\in D$, 则对于 $\forall A\in D$, $A$ 与 $A_0$ 之间均可用 $D$ 内无穷多条分段光滑的曲线连接. 如果我们写出其中任意一条曲线 $\Gamma$ (以 $A_0$ 为起点, 以 $A$ 为终点) 的参数方程, 然后计算可知
\[
  \int_\Gamma u_x'\mathrm dx+u_y'\mathrm dy+u_z'\mathrm dz=u(A)-u(A_0).
\]
换句话说, $u_x'\mathrm dx+u_y'\mathrm dy+u_z'\mathrm dz$ 在 $D$ 内沿任何光滑曲线上的积分, 只依赖于该曲线的起点与终点, 而与曲线的选取无关. 对于这种现象, 我们引入以下概念。

\noindent\textbf{定义 16.7.1}\quad 设 $D\subset\mathbb R^3$ 为一个区域, 函数 $P(x,y,z),Q(x,y,z),R(x,y,z)$ 在 $D$ 内连续. 如果对于 $\forall A_0,A\in D$ 及 $D$ 内任意一条以 $A_0$ 为起点, $A$ 为终点的分段光滑的曲线 $\Gamma$, 曲线积分
\[
  \int_\Gamma P\mathrm dx+Q\mathrm dy+R\mathrm dz
\]
的值只与 $A_0,A$ 有关, 而与曲线 $\Gamma$ 的选取无关, 则称曲线积分 $\int_\Gamma P\mathrm dx+Q\mathrm dy+R\mathrm dz$ 在 $D$ 内与路径无关。

从第二型曲线积分的物理背景可以看出, 该曲线积分与路径无关是指 $D$ 内力场 $\boldsymbol F=(P,Q,R)$ 将单位质量的粒子从点 $A_0$ 移动到点 $A$ 所做的功不依赖于粒子所走的路线, 而只依赖粒子运动的起点与终点. 从直观上看, 重力做功应具有该性质, 因为重力对质点所做的功只依赖于空间这两点之间的垂直距离。

现在, 我们来研究对于区域 $D$ 内给定的三个函数 $P,Q,R$, 何时曲线积分 $\int_\Gamma P\mathrm dx+Q\mathrm dy+R\mathrm dz$ 在 $D$ 内与路径无关. 设 $P(x,y,z),Q(x,y,z),R(x,y,z)$ 是空间 $\mathbb R^3$ 中区域 $D$ 内的三个函数, 若存在 $D$ 内的可微函数 $u=u(x,y,z)$, 使得
\[
  \mathrm du=P\mathrm dx+Q\mathrm dy+R\mathrm dz,
\]
则称 $u$ 为 $P\mathrm dx+Q\mathrm dy+R\mathrm dz$ 的一个原函数. 因此, $P\mathrm dx+Q\mathrm dy+R\mathrm dz$ 有原函数 $u$ 的等价定义是: 存在可微函数 $u$, 使得 $P\mathrm dx+Q\mathrm dy+R\mathrm dz$ 是 $u$ 的全微分。

在本节中, 我们将主要证明以下两个定理。

\noindent\textbf{定理 16.7.1}\quad 设 $D$ 是 $\mathbb R^3$ 中的区域, 函数 $P(x,y,z),Q(x,y,z),R(x,y,z)$ 在 $D$ 内连续, 则以下三个命题等价:

(1) 曲线积分 $\int_\Gamma P\mathrm dx+Q\mathrm dy+R\mathrm dz$ 在 $D$ 内与路径无关;

(2) 对 $D$ 内任何分段光滑的约当曲线 $\Gamma$, 有
\[
  \oint_\Gamma P\mathrm dx+Q\mathrm dy+R\mathrm dz=0;
\]

(3) $P\mathrm dx+Q\mathrm dy+R\mathrm dz$ 在 $D$ 内存在原函数。

\noindent\textbf{证明}\quad 先证 $(1)\Rightarrow(2)$. 设 $\Gamma$ 是 $D$ 内任一条分段光滑的约当曲线. 在 $\Gamma$ 上取定两点 $A_0,A$, 则 $\Gamma$ 构成了从点 $A_0$ 到点 $A$ 的两条曲线 $\Gamma_1,\Gamma_2$, 其中 $\Gamma_1$ 取曲线 $\Gamma$ 的正向, $\Gamma_2$ 则取 $\Gamma$ 的负向 (如图 16.7.1). 由 (1) 我们有
\[
  0=\int_{\Gamma_1}P\mathrm dx+Q\mathrm dy+R\mathrm dz
  -\int_{\Gamma_2}P\mathrm dx+Q\mathrm dy+R\mathrm dz
\]
\[
  =\int_{\Gamma_1+\Gamma_2^-}P\mathrm dx+Q\mathrm dy+R\mathrm dz
  =\int_\Gamma P\mathrm dx+Q\mathrm dy+R\mathrm dz,
\]
因此 (2) 成立。

再证 $(2)\Rightarrow(1)$. 对于 $\forall A_0,A\in D$, 任取两条以 $A_0$ 为起点, $A$ 为终点的光滑曲线 $\Gamma_1,\Gamma_2$. 若 $\Gamma_1\cup\Gamma_2^-$ 组成一条约当曲线, 令 $\Gamma=\Gamma_1\cup\Gamma_2^-$, 则我们有
\[
  \int_{\Gamma_1}P\mathrm dx+Q\mathrm dy+R\mathrm dz
  -\int_{\Gamma_2}P\mathrm dx+Q\mathrm dy+R\mathrm dz
\]
\[
  =\int_{\Gamma_1\cup\Gamma_2^-}P\mathrm dx+Q\mathrm dy+R\mathrm dz
  =\int_\Gamma P\mathrm dx+Q\mathrm dy+R\mathrm dz=0.
\]
当 $\Gamma_1\cup\Gamma_2^-$ 不是约当曲线时, 在 $D$ 内选取一条从点 $A_0$ 到点 $A$ 的光滑曲线 $\Gamma_3$, 使得 $\Gamma_1\cup\Gamma_3^-,\Gamma_2\cup\Gamma_3^-$ 均是约当曲线, 如图 16.7.2. 由上面所证有
\[
  \int_{\Gamma_1}P\mathrm dx+Q\mathrm dy+R\mathrm dz
  =\int_{\Gamma_3}P\mathrm dx+Q\mathrm dy+R\mathrm dz
  =\int_{\Gamma_2}P\mathrm dx+Q\mathrm dy+R\mathrm dz,
\]
因此 (1) 成立。

现证 $(1)\Rightarrow(3)$. 取定 $A_0=(x_0,y_0,z_0)\in D$, 任取 $A=(x,y,z)\in D$, 则可记
\[
  u(x,y,z)=\int_{\Gamma_0}P(\xi,\eta,\zeta)\mathrm d\xi
  +Q(\xi,\eta,\zeta)\mathrm d\eta
  +R(\xi,\eta,\zeta)\mathrm d\zeta,
\]
其中 $\Gamma_0$ 为 $D$ 中一条从点 $A_0$ 到点 $A$ 的固定的光滑曲线。

由于 $A\in D$, 因此 $A$ 是 $D$ 的内点. 对于充分小的 $\Delta x$, 有 $A_{\Delta x}=(x+\Delta x,y,z)\in D$, 于是线段 $\overline{AA_{\Delta x}}\subset D$. 记 $\Gamma_{\Delta x}=\Gamma_0\cup\overline{AA_{\Delta x}}$, 如图 16.7.3, 则
\[
  u(x+\Delta x,y,z)-u(x,y,z)
  =
  \int_{\Gamma_{\Delta x}}P\mathrm d\xi+Q\mathrm d\eta+R\mathrm d\zeta
  -\int_{\Gamma_0}P\mathrm d\xi+Q\mathrm d\eta+R\mathrm d\zeta
\]
\[
  =\int_{\overline{AA_{\Delta x}}}P\mathrm d\xi+Q\mathrm d\eta+R\mathrm d\zeta
  =\int_x^{x+\Delta x}P(\xi,\eta,\zeta)\mathrm d\xi
  =P(x+\theta\Delta x,y,z)\Delta x\quad(0<\theta<1).
\]
因此
\[
  \frac{\partial u}{\partial x}
  =\lim_{\Delta x\to0}\frac{u(x+\Delta x,y,z)-u(x,y,z)}{\Delta x}
  =\lim_{\Delta x\to0}P(x+\theta\Delta x,y,z)=P(x,y,z).
\]
同理可证 $\dfrac{\partial u}{\partial y}=Q$ 及 $\dfrac{\partial u}{\partial z}=R$。

由 $P,Q,R$ 的连续性知, $\mathrm du=P\mathrm dx+Q\mathrm dy+R\mathrm dz$ 在 $D$ 内成立, 因此 $u(x,y,z)$ 是 $P\mathrm dx+Q\mathrm dy+R\mathrm dz$ 的一个原函数。

最后证 $(3)\Rightarrow(1)$. 设 $u(x,y,z)$ 满足 $\mathrm du=P\mathrm dx+Q\mathrm dy+R\mathrm dz$, 再设 $A_0=(x_0,y_0,z_0),A_1=(x_1,y_1,z_1)\in D$ 是任意两点, 任取一条从点 $A_0$ 到点 $A$ 的分段光滑曲线
\[
  \Gamma:\begin{cases}
    x=x(t),\\
    y=y(t),\\
    z=z(t),
  \end{cases}\quad t\in[\alpha,\beta],
\]
且 $(x(\alpha),y(\alpha),z(\alpha))=(x_0,y_0,z_0),(x(\beta),y(\beta),z(\beta))=(x_1,y_1,z_1)$, 则
\[
  \int_\Gamma P\mathrm dx+Q\mathrm dy+R\mathrm dz
\]
\[
  =\int_\alpha^\beta
  [P(x(t),y(t),z(t))x'(t)+Q(x(t),y(t),z(t))y'(t)
  +R(x(t),y(t),z(t))z'(t)]\mathrm dt
\]
\[
  =\int_\alpha^\beta\frac{\mathrm du(x(t),y(t),z(t))}{\mathrm dt}\mathrm dt
  =u(x(\beta),y(\beta),z(\beta))-u(x(\alpha),y(\alpha),z(\alpha))
\]
\[
  =u(x_1,y_1,z_1)-u(x_0,y_0,z_0).
\]
这说明了曲线积分 $\int_\Gamma P\mathrm dx+Q\mathrm dy+R\mathrm dz$ 在 $D$ 内与路径无关. 证毕。

当曲线积分 $\int_\Gamma P\mathrm dx+Q\mathrm dy+R\mathrm dz$ 在 $D$ 内与路径无关时, 对于 $A_0=(x_0,y_0,z_0),A_1=(x_1,y_1,z_1)\in D$, 我们用
\[
  \int_{(x_0,y_0,z_0)}^{(x_1,y_1,z_1)}P\mathrm dx+Q\mathrm dy+R\mathrm dz
\]
来表示 $D$ 内从点 $A_0$ 到点 $A_1$ 的任一条路径的曲线积分。

为了得到进一步的判定曲线积分与路径无关的准则, 我们需要对空间区域作一些限制. 设 $D\subset\mathbb R^3$ 是一个区域, 如果 $D$ 内任何一条约当曲线 $\Gamma$ 总能在 $D$ 内连续地缩成 $D$ 内的一点, 则称 $D$ 是一个单连通区域. 因此对于 $\mathbb R^3$ 中的单连通区域 $D$, $D$ 内部可以有“洞”, 如同心球面围成的区域就是单连通的. 而对于形如轮胎的环面所围区域就不是单连通的了. 空间的单连通区域 $D$ 的一个特点是: 对于 $D$ 中任何一条分段光滑的约当曲线 $\Gamma$, 我们总可以在 $D$ 内找一个分片光滑的双侧曲面 $S$, 使得 $\Gamma=\partial S$. 上述事实的证明超出了数学分析课程的范围, 在这里我们就不介绍了. 但是在下面定理的证明中我们将不加证明地利用该事实。

\noindent\textbf{定理 16.7.2}\quad 设 $D\subset\mathbb R^3$ 是单连通区域, 函数 $P(x,y,z),Q(x,y,z),R(x,y,z)$ 在 $D$ 内具有连续偏导数, 则曲线积分 $\int_\Gamma P\mathrm dx+Q\mathrm dy+R\mathrm dz$ 在 $D$ 内与路径无关的充分必要条件是: 在 $D$ 内有
\[
  \frac{\partial R}{\partial y}\equiv\frac{\partial Q}{\partial z},
  \quad
  \frac{\partial P}{\partial z}\equiv\frac{\partial R}{\partial x},
  \quad
  \frac{\partial Q}{\partial x}\equiv\frac{\partial P}{\partial y},
\]
即向量
\[
  \begin{vmatrix}
    \boldsymbol i&\boldsymbol j&\boldsymbol k\\
    \frac{\partial}{\partial x}&\frac{\partial}{\partial y}&\frac{\partial}{\partial z}\\
    P&Q&R
  \end{vmatrix}
  \equiv0.
\]

\noindent\textbf{证明}\quad 必要性\quad 设 $\int_\Gamma P\mathrm dx+Q\mathrm dy+R\mathrm dz$ 在 $D$ 内与路径无关, 则由定理 16.7.1 知, 存在定义在 $D$ 内的可微函数 $u(x,y,z)$, 使得
\[
  P\mathrm dx+Q\mathrm dy+R\mathrm dz=u_x'\mathrm dx+u_y'\mathrm dy+u_z'\mathrm dz.
\]
因此, 在 $D$ 内恒有
\[
  \frac{\partial R}{\partial y}
  =\frac{\partial^2u}{\partial y\partial z}
  =\frac{\partial^2u}{\partial z\partial y}
  =\frac{\partial Q}{\partial z},
  \quad
  \frac{\partial P}{\partial z}
  =\frac{\partial^2u}{\partial z\partial x}
  =\frac{\partial^2u}{\partial x\partial z}
  =\frac{\partial R}{\partial x},
\]
\[
  \frac{\partial Q}{\partial x}
  =\frac{\partial^2u}{\partial x\partial y}
  =\frac{\partial^2u}{\partial y\partial x}
  =\frac{\partial P}{\partial y},
\]
即必要性成立。

充分性\quad 任取 $D$ 内一条分段光滑的约当曲线 $\Gamma$, 作 $D$ 内一个分片光滑的双侧曲面 $S$, 使得 $\partial S=\Gamma$, 且选取 $S$ 的侧使得 $\Gamma$ 定向为正向. 由斯托克斯公式我们有
\[
  \oint_\Gamma P\mathrm dx+Q\mathrm dy+R\mathrm dz
  =
  \iint_S
  \begin{vmatrix}
    \mathrm dy\mathrm dz&\mathrm dz\mathrm dx&\mathrm dx\mathrm dy\\
    \frac{\partial}{\partial x}&\frac{\partial}{\partial y}&\frac{\partial}{\partial z}\\
    P&Q&R
  \end{vmatrix}
  =0.
\]
由定理 16.7.1 知, $\int_\Gamma P\mathrm dx+Q\mathrm dy+R\mathrm dz$ 在 $D$ 内与路径无关. 证毕。

由于我们在很多时候要讨论平面区域内的曲线积分与路径的关系, 我们不加证明地给出下面定理。

\noindent\textbf{定理 16.7.3}\quad 设 $D\subset\mathbb R^2$ 是一个区域, 函数 $P(x,y),Q(x,y)$ 在 $D$ 内连续, 则以下断言等价:

(1) 曲线积分 $\int_\Gamma P\mathrm dx+Q\mathrm dy$ 在 $D$ 内与路径无关;

(2) 对 $D$ 内任何分段光滑的约当曲线 $\Gamma$, 有 $\oint_\Gamma P\mathrm dx+Q\mathrm dy=0$;

(3) 在 $D$ 内存在可微函数 $u(x,y)$, 使得 $\mathrm du=P\mathrm dx+Q\mathrm dy$。

当 $D$ 是单连通区域且 $P,Q$ 具有连续偏导数时, 则下述断言与上述三个断言等价:
\[
  (4)\quad \text{在 }D\text{ 内成立 }\frac{\partial Q}{\partial x}\equiv\frac{\partial P}{\partial y}.
\]
此定理的证明与定理 16.7.1 和 16.7.2 证明方法相同, 不过在证明 (4) 时, 要用到格林公式。

\noindent\textbf{例 16.7.1}\quad 记 $E=\{(x,y):x\in[0,1],y=0\}$, 证明曲线积分
\[
  \int_\Gamma\frac{x\mathrm dy-y\mathrm dx}{x^2+y^2}
  -\frac{(x-1)\mathrm dy-y\mathrm dx}{(x-1)^2+y^2}
\]
在 $D=\mathbb R^2\setminus E$ 内与路径无关。

\noindent\textbf{证明}\quad 设 $\Gamma_0\subset D$ 是任一条分段光滑的约当曲线, 记
\[
  I=\int_{\Gamma_0}\frac{x\mathrm dy-y\mathrm dx}{x^2+y^2}
  -\frac{(x-1)\mathrm dy-y\mathrm dx}{(x-1)^2+y^2},
\]
则
\[
  I=\int_{\Gamma_0}\frac{x\mathrm dy-y\mathrm dx}{x^2+y^2}
  -\int_{\Gamma_0}\frac{(x-1)\mathrm dy-y\mathrm dx}{(x-1)^2+y^2}.
\]
由于
\[
  \frac{\partial}{\partial x}\left(\frac{x}{x^2+y^2}\right)
  =\frac{y^2-x^2}{(x^2+y^2)^2}
  =\frac{\partial}{\partial y}\left(\frac{-y}{x^2+y^2}\right)
\]
及
\[
  \frac{\partial}{\partial x}\left(\frac{x-1}{(x-1)^2+y^2}\right)
  =\frac{y^2-(x-1)^2}{[(x-1)^2+y^2]^2}
  =\frac{\partial}{\partial y}\left(\frac{-y}{(x-1)^2+y^2}\right),
\]
因此, 对任意的光滑约当曲线 $\Gamma_0\subset D$, 设 $\Gamma_0$ 所围区域为 $D_0$, 若 $D_0\cap E=\varnothing$, 则由定理 16.7.3 知 $I=0$. 当 $E\subset D_0$ 时, 点 $(0,0)$ 及 $(1,0)$ 均在 $D_0$ 内, 从而由例 16.5.2 知
\[
  \int_{\Gamma_0}\frac{x\mathrm dy-y\mathrm dx}{x^2+y^2}=2\pi,
  \quad
  \int_{\Gamma_0}\frac{(x-1)\mathrm dy-y\mathrm dx}{(x-1)^2+y^2}=2\pi.
\]
因此同样有 $I=0$. 由定理 16.7.3 知, 曲线积分
\[
  \int_\Gamma\frac{x\mathrm dy-y\mathrm dx}{x^2+y^2}
  -\frac{(x-1)\mathrm dy-y\mathrm dx}{(x-1)^2+y^2}
\]
在 $D$ 内与路径无关。

\noindent\textbf{注}\quad 显然, 若在 $\mathbb R^2\setminus\{(0,0),(0,1)\}$ 内考虑曲线积分
\[
  \int_\Gamma\frac{x\mathrm dy-y\mathrm dx}{x^2+y^2}
  -\frac{(x-1)\mathrm dy-y\mathrm dx}{(x-1)^2+y^2},
\]
则它在 $D$ 内与路径有关. 另外, 我们容易看出, 在 $\mathbb R^2\setminus\{(-\infty,0]\cup[1,+\infty)\}$ 内该曲线积分与路径无关. 以后从复变函数课程中, 我们还将讨论上述的积分问题, 在那里我们可以更加清楚地知道这些曲线积分与路径无关的本质。

\noindent\textbf{例 16.7.2}\quad 设 $f(r)$ 是定义在 $(0,+\infty)$ 内的一元连续函数, 其中 $r=\sqrt{x^2+y^2+z^2}$. 证明: 当区域 $D\subset\mathbb R^3$ 不含原点时, 曲线积分
\[
  \int_\Gamma f(r)(x\mathrm dx+y\mathrm dy+z\mathrm dz)
\]
在 $D$ 内与路径无关。

\noindent\textbf{证明}\quad 由于 $rf(r)$ 连续, 从而它存在原函数 $u(r)$. 取 $u(\sqrt{x^2+y^2+z^2})$, 则当 $D$ 不含原点时, 有
\[
  \mathrm du
  =\frac{\partial u}{\partial r}\cdot\frac{\partial r}{\partial x}\mathrm dx
  +\frac{\partial u}{\partial r}\cdot\frac{\partial r}{\partial y}\mathrm dy
  +\frac{\partial u}{\partial r}\cdot\frac{\partial r}{\partial z}\mathrm dz
\]
\[
  =f(r)r\left(\frac{x}{r}\mathrm dx+\frac{y}{r}\mathrm dy+\frac{z}{r}\mathrm dz\right)
  =f(r)(x\mathrm dx+y\mathrm dy+z\mathrm dz).
\]
由定理 16.7.1 知 $\int_\Gamma f(r)(x\mathrm dx+y\mathrm dy+z\mathrm dz)$ 在 $D$ 内与路径无关。

\noindent\textbf{注}\quad 选择不同的 $f(r)$, 我们可以得到许多在 $\mathbb R^3$ 中不含原点的区域内的函数 $P,Q,R$, 使得 $\int_\Gamma P\mathrm dx+Q\mathrm dy+R\mathrm dz$ 在 $D$ 内与路径无关。

最后我们介绍一下如何来求原函数, 即给定微分 $P\mathrm dx+Q\mathrm dy+R\mathrm dz$, 当 $\int_\Gamma P\mathrm dx+Q\mathrm dy+R\mathrm dz$ 在区域 $D$ 内与路径无关时, 如何找到在 $D$ 内的函数 $u(x,y,z)$, 使得 $\mathrm du=P\mathrm dx+Q\mathrm dy+R\mathrm dz$。

第一种方法是不定积分法: 由于
\[
  \frac{\partial u}{\partial x}=P,\quad
  \frac{\partial u}{\partial y}=Q,\quad
  \frac{\partial u}{\partial z}=R,
\]
将 $y,z$ 视为常数, 对 $\dfrac{\partial u}{\partial x}=P$ 关于 $x$ 求不定积分得
\[
  u(x,y,z)=\int P\mathrm dx+\varphi(y,z),
\]
其中 $\varphi(y,z)$ 是与 $x$ 无关的 $(y,z)$ 的待定函数. 因此我们有
\[
  \frac{\partial\varphi}{\partial y}
  =Q-\frac{\partial}{\partial y}\int P\mathrm dx.
\]
在上面等式两边对 $y$ 求不定积分得
\[
  \varphi(y,z)=\int\left(Q-\frac{\partial}{\partial y}\int P\mathrm dx\right)\mathrm dy+\psi(z),
\]
其中 $\psi(z)$ 是与 $y$ 无关的 $z$ 的待定函数. 因此我们有
\[
  \frac{\partial\psi}{\partial z}
  =
  R-\frac{\partial}{\partial z}\int P\mathrm dx
  -\frac{\partial}{\partial z}
  \left(\int\left(Q-\frac{\partial}{\partial y}\int P\mathrm dx\right)\mathrm dy\right).
\]
在上面等式两边对 $z$ 求不定积分得
\[
  \psi(z)=\int\left[
  R-\frac{\partial}{\partial z}\int P\mathrm dx
  -\frac{\partial}{\partial z}
  \left(\int\left(Q-\frac{\partial}{\partial y}\int P\mathrm dx\right)\mathrm dy\right)
  \right]\mathrm dz.
\]
最后我们有
\[
  u(x,y,z)=\int P\mathrm dx
  +\int\left(Q-\frac{\partial}{\partial y}\int P\mathrm dx\right)\mathrm dy
\]
\[
  \quad+\int\left[
  R-\frac{\partial}{\partial z}\int P\mathrm dx
  -\frac{\partial}{\partial z}
  \left(\int\left(Q-\frac{\partial}{\partial y}\int P\mathrm dx\right)\mathrm dy\right)
  \right]\mathrm dz.
\]

第二种方法是曲线积分法: 给定 $D$ 中一点 $A_0=(x_0,y_0,z_0)$, 对于任意一点 $A=(x,y,z)\in D$, 在 $D$ 内选择一些连接 $A_0,A$ 且易于求曲线积分的路径, 然后计算曲线积分即得。

\noindent\textbf{例 16.7.3}\quad 在 $\mathbb R^3$ 内求 $u(x,y,z)$, 使得
\[
  \mathrm du(x,y,z)=(yze^{xyz}+3x^2)\mathrm dx+(xze^{xyz}+\sin y)\mathrm dy+(xye^{xyz}+2z)\mathrm dz.
\]

\noindent\textbf{解}\quad 记
\[
  P(x,y,z)=yze^{xyz}+3x^2,\quad Q(x,y,z)=xze^{xyz}+\sin y,
\]
\[
  R(x,y,z)=xye^{xyz}+2z.
\]
容易验证它们处处成立
\[
  \frac{\partial R}{\partial y}=\frac{\partial Q}{\partial z},\quad
  \frac{\partial P}{\partial z}=\frac{\partial R}{\partial x},\quad
  \frac{\partial Q}{\partial x}=\frac{\partial P}{\partial y},
\]
因此 $\int_\Gamma P\mathrm dx+Q\mathrm dy+R\mathrm dz$ 在 $\mathbb R^3$ 内与路径无关。

下面我们用两种方法来求 $u(x,y,z)$。

方法 1\quad 由 $\dfrac{\partial u}{\partial x}=yze^{xyz}+3x^2$, 两边对 $x$ 求不定积分得
\begin{equation}
  u(x,y,z)=e^{xyz}+x^3+\varphi(y,z),
  \tag{16.7.1}
\end{equation}
其中 $\varphi(y,z)$ 为与 $x$ 无关的 $(y,z)$ 的待定函数. 上式两边对 $y$ 求偏导数得
\begin{equation}
  Q(x,y,z)=xze^{xyz}+\sin y
  =\frac{\partial u}{\partial y}
  =xze^{xyz}+\frac{\partial\varphi}{\partial y},
  \tag{16.7.2}
\end{equation}
因此 $\dfrac{\partial\varphi}{\partial y}=\sin y$. 该式两边对 $y$ 求不定积分得
\begin{equation}
  \varphi(y,z)=-\cos y+\psi(z),
  \tag{16.7.3}
\end{equation}
其中 $\psi(z)$ 为与 $x,y$ 无关的 $z$ 的待定函数. 将式 (16.7.3) 代入式 (16.7.1), 再对 $z$ 求偏导数得
\[
  R(x,y,z)=xye^{xyz}+2z=\frac{\partial u}{\partial z}=xye^{xyz}+\psi'(z),
\]
从而 $\psi'(z)=2z$. 此式两边求不定积分得 $\psi(z)=z^2+C$, 其中 $C$ 为常数. 因此
\[
  u(x,y,z)=e^{xyz}+x^3-\cos y+z^2+C.
\]

方法 2\quad 设 $\mathbb R^3$ 中任意一点的坐标为 $(\xi,\eta,\zeta)$, 并设 $\Gamma_1,\Gamma_2,\Gamma_3$ 分别是连接 $(0,0,0)$ 与 $(\xi,0,0)$, $(\xi,0,0)$ 与 $(\xi,\eta,0)$, $(\xi,\eta,0)$ 与 $(\xi,\eta,\zeta)$ 的线段, 则 $\Gamma_1\cup\Gamma_2\cup\Gamma_3$ 是连接 $(0,0,0)$ 与 $(\xi,\eta,\zeta)$ 的一条分段光滑曲线。

由于曲线积分与路径无关, 我们得到 $P\mathrm dx+Q\mathrm dy+R\mathrm dz$ 的一个原函数如下:
\[
  u_1(\xi,\eta,\zeta)
  =\int_{(0,0,0)}^{(\xi,\eta,\zeta)}P\mathrm dx+Q\mathrm dy+R\mathrm dz
\]
\[
  =\int_0^\xi P(x,0,0)\mathrm dx
  +\int_0^\eta Q(\xi,y,0)\mathrm dy
  +\int_0^\zeta R(\xi,\eta,z)\mathrm dz
\]
\[
  =\int_0^\xi 3x^2\mathrm dx+\int_0^\eta\sin y\mathrm dy
  +\int_0^\zeta(\xi\eta e^{\xi\eta z}+2z)\mathrm dz
\]
\[
  =\xi^3-\cos\eta+1+e^{\xi\eta\zeta}-1+\zeta^2
  =e^{\xi\eta\zeta}+\xi^3-\cos\eta+\zeta^2.
\]
将 $(\xi,\eta,\zeta)$ 换回 $(x,y,z)$, 并注意到两个原函数可以相差一个常数, 我们有
\[
  u(x,y,z)=e^{xyz}+x^3-\cos y+z^2+C.
\]

\section{场论简介}

场的概念来源于物理学. 在物理学中, 通常将一些量在空间区域的分布称为场. 最常见的场有密度场、温度场、引力场、磁场等. 如果一个场可以用一个函数来刻画, 则称该场为数量场; 如果一个场需要用一个向量函数来刻画, 则称该场为向量场. 例如, 密度场与温度场都是数量场, 而引力场与磁场则都是向量场。

有些场只与空间的位置有关, 而与时间的变化无关, 这些场称为稳定场, 如地球的引力场就是稳定场. 而温度场等随着空间位置及时间的变化而变化, 这样的场则称为不稳定场。

\subsection{数量场的梯度}

前面我们已经介绍了梯度的概念, 在这里我们来回顾一下. 设 $u=u(x,y,z)$ 是区域 $D\subset\mathbb R^3$ 内的一个数量场 (我们这样来表示数量场, 实际上是建立了一个直角坐标系, 从而数量场在该坐标系下可由 $u=u(x,y,z)$ 来表示), 则当 $u$ 是可微函数时, $u$ 的梯度为
\[
  \operatorname{grad}u=
  \left(\frac{\partial u}{\partial x},\frac{\partial u}{\partial y},\frac{\partial u}{\partial z}\right).
\]
设 $\boldsymbol x_0=(x_0,y_0,z)\in D$ 且 $\boldsymbol v=(\cos\alpha,\cos\beta,\cos\gamma)$ 为一个单位向量, 则过 $A_0$ 沿该单位向量的方向导数为
\[
  \left.\frac{\partial f(\boldsymbol x_0)}{\partial\boldsymbol v}\right|_{\boldsymbol x_0}
  =
  \left.\left(\frac{\partial u}{\partial x}\cos\alpha+\frac{\partial u}{\partial y}\cos\beta+\frac{\partial u}{\partial z}\cos\gamma\right)\right|_{\boldsymbol x_0}
  =
  \left.\operatorname{grad}u\right|_{\boldsymbol x_0}\cdot\boldsymbol v.
\]
因此 $\left.\operatorname{grad}u\right|_{\boldsymbol x_0}$ 与 $\boldsymbol v$ 同向时, 在 $\boldsymbol x_0$ 处沿 $\boldsymbol v$ 的方向导数达到最大. 换句话说, 一个数量场的梯度是一个向量; 它是数量场改变最快的方向, 其向量的模长为数量场沿该方向的变化率. 由此定义的梯度显然与表示该数量场的坐标系的选取无关。

梯度与数量场的等位面密切相关. 所谓等位面是指区域 $D$ 内满足方程 $u(x,y,z)=C(C$ 为常数) 的根的集合. 当 $u(x,y,z)$ 具有较好的性质时 (如是 $C^1$ 的), 数量场的等位面一般是一个曲面. 显然, 任何两个不同的等位面均不相交, 且 $D$ 内任一点都在某一等位面上. 等位面可能是空集、一个点或者一个曲面. 当等位面是一曲面时, 前面我们已经知道它的法向量为
\[
  \left(\frac{\partial u}{\partial x},\frac{\partial u}{\partial y},\frac{\partial u}{\partial z}\right).
\]
由此, 等位面在点 $\boldsymbol x\in D$ 处的法向量即为 $u$ 在 $\boldsymbol x$ 处的梯度 $\left.\operatorname{grad}u\right|_{\boldsymbol x}$。

\subsection{向量场的向量线}

设 $D\subset\mathbb R^3$ 为一个区域, 在 $D$ 内有一个向量场
\[
  \boldsymbol F(x,y,z)=(P(x,y,z),Q(x,y,z),R(x,y,z)),
\]
其中 $P,Q,R$ 在 $D$ 内均具有连续偏导数. 现在设 $\Gamma$ 为 $D$ 内的一条光滑曲线, 其方程为
\[
  \begin{cases}
    x=x(t),\\
    y=y(t),\\
    z=z(t),
  \end{cases}\quad t\in I,
\]
则该曲线上任一点处的切向量为 $(x'(t),y'(t),z'(t))$. 若对于 $\forall t\in I$, 该切向量与向量 $\boldsymbol F(x(t),y(t),z(t))$ 平行, 则曲线 $\Gamma$ 称为 $\boldsymbol F$ 的一条向量线. 由于 $(x'(t),y'(t),z'(t))$ 平行于 $(P,Q,R)$, 我们有
\[
  \frac{x'(t)\mathrm dt}{P}=\frac{y'(t)\mathrm dt}{Q}=\frac{z'(t)\mathrm dt}{R},
\]
即 $\Gamma$ 上的点必须满足
\[
  \frac{\mathrm dx}{P}=\frac{\mathrm dy}{Q}=\frac{\mathrm dz}{R}.
\]
因此我们也称满足上述方程的曲线为 $\boldsymbol F$ 的向量线. 对于 $D$ 内的 $C^1$ 向量函数来说, 过 $D$ 内每一点有且只有一条向量线, 并且不同的向量线互不相交, 所有向量线充满区域 $D$. 关于向量线理论在后续课程中还有介绍。

\noindent\textbf{例 16.8.1}\quad 设向量场 $\boldsymbol F(x,y,z)=(x,y,z)$, 求 $\boldsymbol F$ 在 $\mathbb R^3\setminus\{(0,0,0)\}$ 的向量线。

\noindent\textbf{解}\quad 设 $(x_0,y_0,z_0)\in\mathbb R^3$, 且 $(x_0,y_0,z_0)\ne(0,0,0)$, 不妨设 $x_0\ne0$. 下面我们来求过点 $(x_0,y_0,z_0)$ 的向量线. 由
\[
  \frac{\mathrm dx}{x}=\frac{\mathrm dy}{y},\quad
  \frac{\mathrm dx}{x}=\frac{\mathrm dz}{z},
\]
知 $y=c_1x,z=c_2x$, 其中 $c_1,c_2$ 为待定常数. 再由该向量线过点 $(x_0,y_0,z_0)$ 知 $c_1=\dfrac{y_0}{x_0},c_2=\dfrac{z_0}{x_0}$ 即
\[
  \begin{cases}
    y=\dfrac{y_0}{x_0}x,\\
    z=\dfrac{z_0}{x_0}x,
  \end{cases}
\]
此即射线 $\dfrac{x}{x_0}=\dfrac{y}{y_0}=\dfrac{z}{z_0}$. 从几何上看, 也可以简单地求出该向量线。

\subsection{向量场的散度}

设 $\boldsymbol F=(P(x,y,z),Q(x,y,z),R(x,y,z))$ 是空间区域 $D$ 内的一个 $C^1$ 向量场 (函数), 则 $\boldsymbol F$ 的散度是一个数量场, 其定义为
\[
  \operatorname{div}\boldsymbol F=
  \frac{\partial P}{\partial x}+\frac{\partial Q}{\partial y}+\frac{\partial R}{\partial z}.
\]
上述定义的散度实际上与坐标系的选取无关, 对此我们可以用高斯公式来证明之. 事实上, 将 $\boldsymbol F$ 看做一个流速场. 任给 $D$ 中一点 $\boldsymbol x_0$, 以 $\boldsymbol x_0$ 为心, $r$ 为半径作一小球体 $V_r$, 使得 $V_r\subset D$, 再取 $\partial V_r$ 外侧的单位法向量为 $\boldsymbol n$, 则
\[
  \iint_{\partial V_r}\boldsymbol F\cdot\boldsymbol n\mathrm dS
  =
  \iiint_{V_r}\operatorname{div}\boldsymbol F\mathrm dV.
\]
由第二型曲面积分的几何意义, 上述等式的左边是该流速场单位时间从内到外通过 $\partial V_r$ 的流量, 由此它不依赖于坐标系的选取. 由重积分的第一中值定理, 有
\[
  \operatorname{div}\boldsymbol F(\boldsymbol x_0)
  =
  \lim_{r\to0+0}
  \frac{\iint_{\partial V_r}\boldsymbol F\cdot\boldsymbol n\mathrm dS}{V_r},
\]
其中我们仍然用 $V_r$ 表示 $V_r$ 的体积. 由此说明散度的定义也不依赖于坐标系的选取, 是一个向量场的本质属性。

对于在空间区域 $D$ 内具有连续偏导数的向量函数 $\boldsymbol F(x,y,z)$, 若它的散度 $\operatorname{div}\boldsymbol F$ 在 $\boldsymbol x_0=(x_0,y_0,z_0)\in D$ 处不为零, 则当 $\operatorname{div}\boldsymbol F(\boldsymbol x_0)>0$ 时, 对于任何一个包含 $\boldsymbol x_0$ 的充分小光滑闭曲面 $S$ 所围的区域 $V$, 有
\[
  \iiint_V\operatorname{div}\boldsymbol F\mathrm dV>0.
\]
由高斯公式知, 此处从内向外流出的流量为正. 因此 $\boldsymbol x_0$ 也称为该向量场的源. 反之, 若 $\operatorname{div}\boldsymbol F(\boldsymbol x_0)<0$, 则对任何内部含有 $\boldsymbol x_0$ 的充分小的闭曲面, 从外面里流进的流量为正. 这时, $\boldsymbol x_0$ 也称该向量场的汇. 当 $\operatorname{div}\boldsymbol F(\boldsymbol x_0)=0$ 时, 称向量场 $\boldsymbol F$ 在 $\boldsymbol x_0$ 处无源。

\subsection{向量场的旋度}

如同上一小节, 我们仍设 $\boldsymbol F(x,y,z)=(P(x,y,z),Q(x,y,z),R(x,y,z))$ 为区域 $D\subset\mathbb R^3$ 内的一个 $C^1$ 向量场. 在斯托克斯公式中, 曾出现如下与 $\boldsymbol F$ 有关的向量
\[
  \begin{vmatrix}
    \boldsymbol i&\boldsymbol j&\boldsymbol k\\
    \frac{\partial}{\partial x}&\frac{\partial}{\partial y}&\frac{\partial}{\partial z}\\
    P&Q&R
  \end{vmatrix}.
\]
我们称此向量为 $\boldsymbol F$ 的旋度, 记为 $\operatorname{rot}\boldsymbol F$。

设 $S$ 是一分片光滑曲面, $\partial S$ 是空间分段光滑闭曲线, 且其定向与 $S$ 的侧的单位法向量 $\boldsymbol n$ 的选定一致, 则斯托克斯公式告诉我们
\[
  \oint_{\partial S}P\mathrm dx+Q\mathrm dy+R\mathrm dz
  =
  \iint_S\operatorname{rot}\boldsymbol F\cdot\boldsymbol n\mathrm dS.
\]
上式左边的线积分通常称为 $\boldsymbol F$ 沿 $\partial S$ 的环量, 它显然可记成 $\oint_{\partial S}\boldsymbol F\mathrm d\boldsymbol s$, 其中 $\mathrm d\boldsymbol s=(\mathrm dx,\mathrm dy,\mathrm dz)$. 由此定义的环量与坐标系的选取无关。

利用这一定义, 我们同样可说明旋度也与坐标系选取无关. 事实上, 任取以 $\boldsymbol x_0\in D$ 为心, 半径 $r$ 充分小的一个圆盘 $S_r\subset D$, 取定 $\partial S_r$ 的方向与 $S_r$ 的一侧的单位法向量 $\boldsymbol n$ 保持一致, 则
\[
  \oint_{\partial S_r}\boldsymbol F\mathrm d\boldsymbol s
  =
  \iint_{S_r}\operatorname{rot}\boldsymbol F\cdot\boldsymbol n\mathrm dS,
\]
从而
\[
  \lim_{r\to0}\operatorname{rot}\boldsymbol F\cdot\boldsymbol n
  =
  \lim_{r\to0}
  \frac{\displaystyle\oint_{\partial S_r}\boldsymbol F\mathrm d\boldsymbol s}{\pi r^2}.
\]
由于 $S_r$ 是任意一个圆盘, 从而 $\boldsymbol n$ 可取到 $\boldsymbol x_0$ 出发的任何方向. 上式说明向量 $\operatorname{rot}\boldsymbol F$ 在任何方向的投影均不依赖于坐标系的选取, 从而 $\operatorname{rot}\boldsymbol F$ 与坐标系的选取无关. 由斯托克斯定理, 旋度与环量是密切相关的. 特别地, 曲线积分 $\int_\Gamma P\mathrm dx+Q\mathrm dy+R\mathrm dz$ 在区域 $D$ 内与路径无关的一个必要条件是 $\boldsymbol F$ 的旋度处处为零. 当 $\boldsymbol F$ 的旋度处处为零时, 我们称 $\boldsymbol F$ 为一个无旋场。

\subsection{一些重要算子}

所谓“算子”是作用在某些特定对象 (如函数、向量函数等) 的一些特别的运算. 人们对它们赋予了一些特殊的记号. 在这里我们主要介绍两个重要的算子。

\noindent\textbf{1. 哈密尔顿算子}

哈密尔顿 (Hamilton) 曾在 $\mathbb R^n$ 中引进了以下算子:
\[
  \nabla=\left(\frac{\partial}{\partial x_1},\frac{\partial}{\partial x_2},\cdots,\frac{\partial}{\partial x_n}\right).
\]
我们称该算子为哈密尔顿算子. 对于区域 $D\subset\mathbb R^n$ 内的可微函数 $f(x_1,x_2,\cdots,x_n)$, 哈密尔顿算子 $\nabla$ 对 $f$ 的作用定义为
\[
  \nabla f=\left(\frac{\partial f}{\partial x_1},\frac{\partial f}{\partial x_2},\cdots,\frac{\partial f}{\partial x_n}\right)=\operatorname{grad}f.
\]

容易验证, 哈密尔顿算子 $\nabla$ 具有如下性质: 设 $c\in\mathbb R$, $f,g$ 是可微函数, 则
\[
  (1)\quad \nabla(cf)=c\nabla f;
\]
\[
  (2)\quad \nabla(f+g)=\nabla f+\nabla g;
\]
\[
  (3)\quad \nabla(fg)=f\nabla g+g\nabla f.
\]

另外, $\mathbb R^3$ 中的向量场 $\boldsymbol F=(P,Q,R)$ 的散度可表示为
\[
  \operatorname{div}\boldsymbol F
  =
  \frac{\partial P}{\partial x}+\frac{\partial Q}{\partial y}+\frac{\partial R}{\partial z}
  =
  \nabla\cdot\boldsymbol F,
\]
而 $\boldsymbol F$ 的旋度表示为
\[
  \operatorname{rot}\boldsymbol F
  =
  \begin{vmatrix}
    \boldsymbol i&\boldsymbol j&\boldsymbol k\\
    \frac{\partial}{\partial x}&\frac{\partial}{\partial y}&\frac{\partial}{\partial z}\\
    P&Q&R
  \end{vmatrix}
  =
  \nabla\times\boldsymbol F.
\]

\noindent\textbf{2. 拉普拉斯算子}

下面我们介绍第二个重要算子, 即拉普拉斯算子. 在 $\mathbb R^n$ 中, 拉普拉斯算子的定义为
\[
  \Delta=\sum_{i=1}^n\frac{\partial^2}{\partial x_i^2}.
\]
对于在区域 $D\subset\mathbb R^n$ 内具有二阶连续偏导数的函数 $f(x_1,x_2,\cdots,x_n)$, 拉普拉斯算子 $\Delta$ 对 $f$ 的作用定义为
\[
  \Delta f=\sum_{i=1}^n\frac{\partial^2f}{\partial x_i^2}.
\]
当 $u=f(x_1,x_2,\cdots,x_n)$ 在 $D$ 内恒满足
\[
  \Delta u=\sum_{i=1}^n\frac{\partial^2f}{\partial x_i^2}=0
\]
时, 称 $u$ 为 $D$ 内的调和函数。

显然一元函数 $f(x)$ 为调和函数的充分必要条件是 $f(x)$ 为线性函数. 对二元函数 $u=f(x,y)$, 由定义, 当 $\Delta u=\dfrac{\partial^2f}{\partial x^2}+\dfrac{\partial^2f}{\partial y^2}=0$ 时, $u=f(x,y)$ 是调和函数. 读者不难验证, $e^x\cos y,2xy$ 均是调和函数. 调和函数具有许多特殊的性质, 在后续课程中将有详细介绍。

\noindent\textbf{例 16.8.2}\quad 设 $D\subset\mathbb R^2$ 是单连通区域, $u(x,y)$ 是 $D$ 内的调和函数. 证明曲线积分
\[
  \int_\Gamma u_y'\mathrm dx-u_x'\mathrm dy
\]
在 $D$ 内与路径无关, 且其原函数为 $D$ 内的调和函数。

\noindent\textbf{证明}\quad 由于 $P\mathrm dx+Q\mathrm dy=u_y'\mathrm dx-u_x'\mathrm dy$ 及
\[
  \frac{\partial Q}{\partial x}-\frac{\partial P}{\partial y}
  =
  \frac{\partial(-u_x')}{\partial x}-\frac{\partial u_y'}{\partial y}
  =
  -\left(\frac{\partial^2u}{\partial x^2}+\frac{\partial^2u}{\partial y^2}\right)
  \equiv0,
\]
再注意到 $D$ 是单连通区域, 根据定理 16.7.3 知, $\int_\Gamma u_y'\mathrm dx-u_x'\mathrm dy$ 在 $D$ 内与路径无关。

设 $v(x,y)$ 为 $u_y'\mathrm dx-u_x'\mathrm dy$ 的一个原函数, 则
\[
  \frac{\partial v}{\partial x}=u_y',
  \quad
  \frac{\partial v}{\partial y}=-u_x'.
\]
因此
\[
  \frac{\partial^2v}{\partial x^2}+\frac{\partial^2v}{\partial y^2}
  =
  \frac{\partial^2u}{\partial y\partial x}
  -
  \frac{\partial^2u}{\partial x\partial y}
  \equiv0,
\]
所以 $v(x,y)$ 是 $D$ 内的调和函数。

\noindent\textbf{注}\quad 设 $u(x,y)$ 是 $D$ 内的调和函数, 在上例中得到的 $D$ 内的调和函数 $v(x,y)$ 称为 $u(x,y)$ 的共轭调和函数。

\section*{习题十六}

1. 设 $\Gamma$ 是空间 $\mathbb R^3$ 中一条光滑的物质曲线, 其密度函数 $\rho(x,y,z)$ 在 $\Gamma$ 上连续, 试求 $\Gamma$ 的质心坐标。

2. 设曲线 $\Gamma$ 是单位圆周: $x^2+y^2=1$, 求下列第一型曲线积分:
\[
  (1)\quad \int_\Gamma x\mathrm ds;\qquad
  (2)\quad \int_\Gamma xy\mathrm ds;\qquad
  (3)\quad \int_\Gamma x^2\mathrm ds;\qquad
  (4)\quad \int_\Gamma |x|\mathrm ds.
\]

3. 设曲线 $\Gamma$ 为球面 $x^2+y^2+z^2=1$ 与平面 $x+y+z=0$ 的交线, 求下列第一型曲线积分:
\[
  (1)\quad \int_\Gamma x\mathrm ds;\qquad
  (2)\quad \int_\Gamma xy\mathrm ds;\qquad
  (3)\quad \int_\Gamma x^2\mathrm ds.
\]

4. 求下列第一型曲线积分:

\end{document}
